\chapter{Análisis de convergencia}
\label{chap:convergencia}

\noindent
En este capítulo se analiza la convergencia de redes neuronales físicamente informadas (PINNs) para el
\emph{problema de Dirichlet elíptico} y la \emph{formulación de pérdida} introducidos en los Capítulos~1 y~2
(véase, en particular, la definición de la pérdida PINN en la Sección~\ref{sec:raissi-loss} y los requisitos
de regularidad en la Sección~\ref{sec:reg-ad-pinns}). El objetivo es conectar rigurosamente tres niveles:
(i) la pérdida empírica computable a partir de puntos de colocación, (ii) la pérdida teórica asociada a normas
$L^2$ respecto de medidas de muestreo, y (iii) el error $h-u^*$ en normas funcionales (primero $C^0$ y luego,
bajo imposición fuerte de la frontera, $H^1$).


\medskip
\noindent\textbf{Observación (Dependencia bibliográfica y aporte del capítulo).}
El paso de \emph{generalización} que conecta la pérdida empírica $\widehat{\mathcal{E}}_m$ con la pérdida teórica $\mathcal{E}$
(Teorema~1 y lemas de control determinista y de cobertura probabilística) sigue la estrategia de cobertura de Shin et al.~\cite{shin2020convergence}.
Un enfoque complementario de estimaciones del error de generalización mediante estabilidad de la EDP y error de entrenamiento es el de Mishra y Molinaro~\cite{mishra2023generalization}. Se cita como contexto y no como fundamento de las demostraciones que siguen.
Las intervenciones se distinguen según su naturaleza: precisiones de
notación, reparaciones de pasos demostrativos y consecuencias derivadas
bajo las hipótesis. La escritura coherente del operador elíptico del
Capítulo~1 no se presenta como aporte original. La elección de pesos
iguales del Teorema~2 evita una comparabilidad no garantizada por
las condiciones generales del artículo. Los argumentos de cobertura
proceden de \cite{shin2020convergence}, mientras que la interpolación
cuantitativa $L^2$--H\"older, la tasa de error y la barrera explícita del
Teorema~3 son desarrollos detallados de consecuencias condicionales,
no afirmaciones sobre optimización numérica ni refutaciones globales del
artículo.
El aporte de esta sección consiste en \emph{desarrollar explícitamente los argumentos de cobertura y sus hipótesis}: (i) mantenemos separadas las medidas
de muestreo $\mu_r$ (interior) y $\mu_b$ (borde) y la posibilidad de muestrear únicamente sobre un subconjunto
$\Gamma\subseteq\partial U$; (ii) desarrollamos el paso determinista ``continuo--discreto'' bajo un evento de $\varepsilon$-cobertura,
incluyendo el cierre interior--borde; y (iii) rastreamos con precisión qué términos dependen de la regularidad Hölder del residuo y
cuáles dependen solo de los datos $f$ y $g$.
\medskip
\noindent\textbf{Alcance lógico de las demostraciones.}
Los Teoremas~1--4 constituyen una cadena \emph{condicional}: la Hipótesis~1
controla el muestreo aleatorio; la Hipótesis~2 exige comparadores con
interpolación exacta y regularidad uniforme, además de la existencia de
minimizadores globales mensurables en el Teorema~2 original; la Sección~\ref{sec:comparadores_aproximados}
permite sustituir ambas exigencias de exactitud por tolerancias cuantitativas; y la Hipótesis~3 añade las condiciones
elípticas necesarias para estabilidad y existencia. El Teorema~4 restringe
adicionalmente la familia de aproximantes a un espacio con condición de
Dirichlet impuesta exactamente. En los experimentos se usan también mallas
deterministas y optimizadores iterativos: tales ensayos no verifican por sí
solos las hipótesis probabilísticas ni la minimización global.

\medskip
\noindent\textbf{Convención de notación.}
Sea $u^*$ la solución exacta del problema elíptico (Capítulo~1) y sea $h$ una aproximación tipo PINN.
Denotamos por
\[
r_h:=\mathcal{L}[h]-f\quad\text{en }U,
\qquad
b_h:=\mathcal{B}[h]-g\quad\text{en }\Gamma
\]
los residuos interior y de frontera (Dirichlet), respectivamente. Para abreviar, escribimos
$\mathcal{E}(h)$ para la \emph{pérdida teórica} y $\widehat{\mathcal{E}}_m(h)$ para la \emph{pérdida empírica}
(en el sentido del Capítulo~2, es decir, $\mathcal{E}(h)=\text{Loss}(h;\lambda,\lambda^{\mathcal{R}})$ y
$\widehat{\mathcal{E}}_m(h)=\text{Loss}_m(h;\lambda,\lambda^{\mathcal{R}})$; véase la Sección~\ref{sec:raissi-loss}), con medidas de muestreo $\mu_r$ en $U$ y $\mu_b$ en $\Gamma$,
y tamaño de muestreo $m=(m_r,m_b)$.

En coherencia con el planteamiento abstracto de \cite{shin2020convergence}, consideraremos un problema de contorno de la forma
\begin{equation}\label{eq:planteamiento_abstracto}
\mathcal{L}[u](x)=f(x)\quad (x\in U),\qquad \mathcal{B}[u](x)=g(x)\quad (x\in\Gamma\subseteq\partial U).
\end{equation}
En particular, $\Gamma$ representa la porción de frontera en la que se impone la condición (y el soporte geométrico de la medida $\mu_b$ usada para muestrear datos de frontera). A lo largo del capítulo, se controlará $\mathcal{E}(h)$ en términos de
$\widehat{\mathcal{E}}_m(h)$ más un término geométrico–probabilístico de cobertura; siguiendo
\cite{shin2020convergence}, se desarrollan aquí las pruebas en detalle y se explicitan los pasos omitidos.

\medskip
\paragraph{Interpretación de los resultados principales.}
El hilo lógico del capítulo puede sintetizarse como la cadena
\[
\widehat{\mathcal{E}}_m(h)\ (\text{discreto}) \;\Rightarrow\; \mathcal{E}(h)\ (\text{continuo})
\;\Rightarrow\; \|r_h\|_{C^0(\overline U)}+\|b_h\|_{C^0(\Gamma)}
\;\Rightarrow\; \|h-u^*\|_{C^0(\overline U)}.
\]
La primera flecha es el contenido del Teorema~1 (control probabilístico vía cobertura).
La segunda se obtiene en el Teorema~2, donde la regularización Hölder permite pasar de control $L^2$ a control uniforme de los residuos.
En la segunda parte del capítulo, desde la Sección~\ref{sec:existencia_dirichlet}, identificamos $\Gamma=\partial U$ para tratar el problema de Dirichlet clásico.
La tercera es un paso puramente elíptico: una estimación de estabilidad transforma el residuo interior y el defecto de contorno en
control del error (estimación de estabilidad de la Sección~\ref{sec:estabilidad_dirichlet}), concluyendo el Teorema~3. El Teorema~4 añade, en el caso de Poisson con Dirichlet fuerte,
convergencia en la norma de energía $H^1$.

\section{Regularidad y precompacidad de redes neuronales (Lema 1)}


\medskip
\noindent\textbf{Observación (Fuente y uso en el capítulo).}
Este lema fija la hipótesis de suavidad mínima para que $\mathcal{L}[h]$ y $\mathcal{B}[h]$ puedan evaluarse punto a punto,
y para que los residuos admitan cotas uniformes de derivadas y seminormas Hölder. El planteamiento es afín al de
\cite[Lema~2.1]{shin2020convergence}, pero aquí se explicita qué requerimos de la activación $\sigma$ y cómo las cotas sobre
parámetros se traducen en equicontinuidad y precompacidad en $C^{k-1}(\overline U)$.
\medskip

Consideremos $U$ un dominio acotado y $\mathcal{H}^{\mathrm{NN}}_{\vec n}$ una clase de redes neuronales de propagación hacia adelante y totalmente conectadas con arquitectura $\vec n=(n_0,\dots,n_L)$, entrada en $\mathbb{R}^{n_0}$ y salida escalar.
Supondremos que la activación $\sigma\in C^{k'}(\mathbb{R})$ para algún $k'>k\ge 1$ y que, para cada
$s\in\{0,\dots,k\}$, la derivada $\sigma^{(s)}$ es acotada y Lipschitz.

\medskip
\noindent\textbf{Enunciado (regularidad y precompacidad).}
Sea $h_\theta\in\mathcal{H}^{\mathrm{NN}}_{\vec n}$, donde $\theta$ denota el conjunto de pesos y sesgos.
\begin{enumerate}
\item Si $\sigma\in C^{k'}$ con $k'>k$, entonces $h_\theta\in C^k(\mathbb{R}^{n_0})$. Además, si los pesos y
sesgos satisfacen $\|\theta\|_2\le M$ (norma euclídea del vector obtenido al concatenar todos los parámetros), entonces existe una
constante $C=C(\vec n,k,M,\sigma)$ tal que $\|D^\beta h_\theta\|_{L^\infty(\overline U)}\le C$ para todo
multiíndice $\beta$ con $|\beta|\le k$.
\item (Precompacidad) Sea $\{h_j\}_{j\ge 1}\subset\mathcal{H}^{\mathrm{NN}}_{\vec n}$ una sucesión con
parámetros uniformemente acotados por $M$ y restringida a $\overline U$. Entonces $\{h_j\}$ es equicontinua
en $\overline U$ y, para cada $\ell\le k-1$, la familia $\{D^\beta h_j:|\beta|\le \ell\}$ es equicontinua.
En particular, existe una subsucesión (no relabelada) tal que $h_j\to h_\infty$ en $C^{k-1}(\overline U)$.
\end{enumerate}

\medskip
\noindent\textbf{Demostración.}

(i) \emph{Regularidad $C^k$ y cotas uniformes: inducción explícita por capas.}
Para hacer completamente verificable el argumento, escribimos la red mediante variables intermedias.
Definimos $z_0(x)=x$ y, para $\ell=1,\dots,L-1$,
\[
a_\ell(x):=W_\ell z_{\ell-1}(x)+b_\ell,\qquad
z_\ell(x):=\sigma(a_\ell(x))\ \ \text{(aplicada componente a componente)}.
\]
Suponemos la última capa lineal, es decir
\[
h_\theta(x)=W_L z_{L-1}(x)+b_L.
\]

\smallskip
\noindent\textit{Paso base ($\ell=0$).}
Se tiene $z_0\in C^\infty$ y
$Dz_0\equiv I$ mientras que $D^\beta z_0\equiv 0$ para todo multiíndice $\beta$ con $|\beta|\ge 2$.
En particular, para cada $s\le k$ existe una constante $M_{0,s}$ (de hecho $M_{0,0}=\sup_{x\in U}\|x\|$, $M_{0,1}=1$ y $M_{0,s}=0$ si
$s\ge 2$) tal que
\[
\sup_{|\beta|=s}\|D^\beta z_0\|_{L^\infty(\overline U)}\le M_{0,s}.
\]

\smallskip
\noindent\textit{Hipótesis inductiva.}
Fijemos $\ell\in\{1,\dots,L-1\}$ y supongamos que $z_{\ell-1}\in C^k(\mathbb{R}^{n_0})$ y que, para cada
$s\in\{0,1,\dots,k\}$, existen constantes $M_{\ell-1,s}$ tales que
\begin{equation}\label{eq:lemma1_ind_hyp}
\sup_{|\beta|=s}\|D^\beta z_{\ell-1}\|_{L^\infty(\overline U)}\le M_{\ell-1,s}.
\end{equation}

\smallskip
\noindent\textit{Paso inductivo: derivadas de la parte afín.}
Como $a_\ell=W_\ell z_{\ell-1}+b_\ell$ y $W_\ell$ es lineal, para todo multiíndice $\beta\neq 0$ se cumple
\[
D^\beta a_\ell(x)=W_\ell\,D^\beta z_{\ell-1}(x),
\]
y por tanto
\begin{equation}\label{eq:lemma1_affine_bound}
\sup_{|\beta|=s}\|D^\beta a_\ell\|_{L^\infty(\overline U)}\le \|W_\ell\|\,M_{\ell-1,s}\qquad (1\le s\le k).
\end{equation}

\smallskip
\noindent\textit{Paso inductivo: derivadas tras la activación.}
Para cada componente $(z_\ell)_j=\sigma((a_\ell)_j)$, la regla de la cadena y su extensión multivariada
(Faà di Bruno) dan, para $s=1$ y $s=2$, las identidades (a nivel tensorial)
\[
D z_\ell=\mathrm{Diag}(\sigma'(a_\ell))\,D a_\ell,
\qquad
D^2 z_\ell=\mathrm{Diag}(\sigma''(a_\ell))\,[Da_\ell,Da_\ell]
\,+\,\mathrm{Diag}(\sigma'(a_\ell))\,D^2 a_\ell,
\]
y, para orden general $s\le k$, que $D^s z_\ell$ es suma finita de términos de la forma
\begin{equation}\label{eq:lemma1_fdb_form}
\sigma^{(t)}(a_\ell)\,\big[ D^{s_1}a_\ell,\dots,D^{s_t}a_\ell\big],
\qquad
1\le t\le s,\quad s_1+\cdots+s_t=s,\quad s_i\ge 1,
\end{equation}
donde los corchetes representan una contracción multilineal universal (dependiente solo de $s,n_0$).

Tomando norma supremo en $\overline U$, usando que $\sigma^{(t)}$ es acotada para $t\le k$ y aplicando
\eqref{eq:lemma1_affine_bound} en cada factor, se obtiene una cota inductiva de la forma
\begin{equation}\label{eq:lemma1_ind_step_bound}
\sup_{|\beta|=s}\|D^\beta z_\ell\|_{L^\infty(\overline U)}
\le
P_{s}\Big(\|W_\ell\|,\;M_{\ell-1,1},\dots,M_{\ell-1,s}\Big)
\qquad (1\le s\le k),
\end{equation}
para ciertos polinomios $P_s$ con coeficientes que dependen solo de $s,n_0$ y de las cotas
$\|\sigma^{(t)}\|_{L^\infty}$, $t\le s$. En particular, si $\|\theta\|\le M$, entonces
$\|W_\ell\|\le M$ para todo $\ell$ y, por iteración en $\ell$, existen constantes $M_{\ell,s}$ (dependientes
solo de $\vec n,k,M,\sigma$) tales que
\[\sup_{|\beta|=s}\|D^\beta z_\ell\|_{L^\infty(\overline U)}\le M_{\ell,s}\quad\text{para todo }s\le k.
\]
Finalmente, como $h_\theta=W_L z_{L-1}+b_L$ es combinación lineal de $z_{L-1}$, el mismo argumento que en
\eqref{eq:lemma1_affine_bound} transfiere las cotas a $D^\beta h_\theta$ para $|\beta|\le k$.

(ii) En este caso las derivadas se obtienen de redes definidas en todo
$\mathbb R^{n_0}$ y las cotas de parámetros dan las mismas estimaciones
sobre la envolvente convexa compacta de $\overline U$. Por ello,
$\{h_j\}$ es equi-Lipschitz en $\overline U$, luego equicontinua. De manera análoga, las cotas uniformes sobre derivadas hasta orden
$k$ implican equicontinuidad de derivadas hasta orden $k-1$ (pues $D^\beta h_j$ es Lipschitz si
$D^{\beta+e_i}h_j$ es acotada). Aplicando el teorema de Arzelà--Ascoli de manera iterada a la familia
$\{D^\beta h_j\}_{|\beta|\le k-1}$ se obtiene una subsucesión convergente en $C^{k-1}(\overline U)$.

\hfill $\square$

\medskip
\noindent\textbf{Comentario técnico.}
En las secciones siguientes no usaremos este lema para deducir convergencia de derivadas a partir de la minimización (eso sería, en general,
demasiado fuerte), sino para justificar que las clases $\mathcal{H}_m$ se eligen dentro de espacios donde los residuos
$\mathcal{L}[h]-f$ y $\mathcal{B}[h]-g$ están bien definidos y admiten controles de regularidad (por ejemplo, seminormas Hölder finitas).
En particular, la precompacidad en $C^{k-1}(\overline U)$ es la base para pasar al límite en argumentos de estabilidad elíptica.


\section{Muestreo y condiciones geométricas de cobertura (Hipótesis 1)}

Fijamos el esquema de muestreo i.i.d. y las hipótesis geométricas de cubrimiento utilizadas en
el argumento probabilístico. El contenido sigue, con notación adaptada,
\cite[Assumption~3.1]{shin2020convergence}, distinguiendo explícitamente la medida volumétrica
$\mu_r$ en $U$ y la medida superficial $\mu_b$ en $\Gamma\subseteq\partial U$. Las cotas
inferiores de masa en bolas que se emplearán más adelante no se incorporan como una hipótesis
independiente: se deducen de las particiones postuladas por Shin et al.; esta consecuencia se
formula en el lema auxiliar de no degeneración local.

\medskip
\noindent\textbf{Hipótesis 1 (muestreo i.i.d. y cubrimiento probabilístico).}
Sea $U\subset\mathbb R^d$ un dominio acotado de clase, al menos, $C^{0,1}$ y sea
$\Gamma\subseteq\partial U$ un subconjunto cerrado. Sea $\mu_r$ una medida de probabilidad
en $U$, absolutamente continua con respecto a la medida de Lebesgue y con densidad $\rho_r$;
sea $\mu_b$ una medida de probabilidad en $\Gamma$, absolutamente continua con respecto a la
medida de Hausdorff de dimensión $d-1$ y con densidad $\rho_b$. Supondremos que las densidades
están soportadas en $U$ y $\Gamma$, respectivamente, y que
\[
\operatorname*{ess\,inf}_{x\in U}\rho_r(x)>0,
\qquad
\operatorname*{ess\,inf}_{y\in\Gamma}\rho_b(y)>0.
\]
Las muestras $X_r=\{x_r^i\}_{i=1}^{m_r}$ y $X_b=\{x_b^j\}_{j=1}^{m_b}$ son i.i.d. con
leyes $\mu_r$ y $\mu_b$, respectivamente, e independientes entre sí.

Además, existen constantes $c_r,C_r,c_b,C_b>0$ y $\varepsilon_0>0$ tales que, para todo
$0<\varepsilon\leq\varepsilon_0$, existen particiones finitas
$\{U_k^{\varepsilon}\}_{k=1}^{K_{r,\varepsilon}}$ de $U$ y
$\{\Gamma_\ell^{\varepsilon}\}_{\ell=1}^{K_{b,\varepsilon}}$ de $\Gamma$ que verifican
\[
\operatorname{diam}(U_k^{\varepsilon})\leq \sqrt d\,\varepsilon,
\qquad
\mu_r(U_k^{\varepsilon})\geq c_r\varepsilon^d,
\]
\[
\operatorname{diam}(\Gamma_\ell^{\varepsilon})\leq \sqrt d\,\varepsilon,
\qquad
\mu_b(\Gamma_\ell^{\varepsilon})\geq c_b\varepsilon^{d-1}.
\]
Asimismo, para todo $x\in U$ y $y\in\Gamma$,
\[
\mu_r\bigl(B(x,\varepsilon)\cap U\bigr)\leq C_r\varepsilon^d,
\qquad
\mu_b\bigl(B(y,\varepsilon)\cap\Gamma\bigr)\leq C_b\varepsilon^{d-1}.
\]
En particular,
\[
K_{r,\varepsilon}\leq(c_r\varepsilon^d)^{-1},
\qquad
K_{b,\varepsilon}\leq(c_b\varepsilon^{d-1})^{-1}.
\]
Cuando $d=1$, siguiendo a Shin et al., se supone que todos los puntos de la frontera están
disponibles y no se requiere un muestreo aleatorio de borde.

\medskip
\noindent
Fijamos además $\alpha\in(0,1]$ y suponemos que, para los candidatos $h$ relevantes,
\[
[r_h]_{\alpha;U}<\infty,
\qquad
[b_h]_{\alpha;\Gamma}<\infty,
\]
lo cual será garantizado por la elección de la clase de redes y por los términos de
regularización $\mathcal R_r$ y $\mathcal R_b$.

\medskip
\noindent\textbf{Lema auxiliar (No degeneración local deducida de las particiones).}
Supóngase $d\geq2$ y que se cumple la Hipótesis~1. Entonces existen constantes
$\underline c_r,\underline c_b>0$ tales que, para todo radio
$0<r\leq2\sqrt d\,\varepsilon_0$,
\[
\mu_r\bigl(B(x,r)\cap U\bigr)\geq \underline c_r r^d,
\qquad x\in U,
\]
y
\[
\mu_b\bigl(B(y,r)\cap\Gamma\bigr)\geq \underline c_b r^{d-1},
\qquad y\in\Gamma.
\]
Puede tomarse
\[
\underline c_r=c_r(2\sqrt d)^{-d},
\qquad
\underline c_b=c_b(2\sqrt d)^{-(d-1)}.
\]
\medskip

\noindent\textit{Demostración.}
Sea $x\in U$ y tome $\varepsilon=r/(2\sqrt d)\leq\varepsilon_0$.
Como $\{U_k^\varepsilon\}$ es una partición, existe una celda
$U_k^\varepsilon$ que contiene a $x$. Para todo $z\in U_k^\varepsilon$,
\[
\|z-x\|\leq\operatorname{diam}(U_k^\varepsilon)
\leq\sqrt d\,\varepsilon=r/2<r.
\]
Luego $U_k^\varepsilon\subset B(x,r)\cap U$ y
\[
\mu_r(B(x,r)\cap U)\geq\mu_r(U_k^\varepsilon)
\geq c_r\varepsilon^d=c_r(2\sqrt d)^{-d}r^d.
\]
El mismo argumento con la partición de $\Gamma$ da
$\mu_b(B(y,r)\cap\Gamma)\geq c_b(2\sqrt d)^{-(d-1)}r^{d-1}$.
Se ha elegido deliberadamente un diámetro $r/2$ para evitar identificar
bolas abiertas con bolas cerradas en el paso de inclusión.
\hfill$\square$

\medskip
\noindent\textbf{Ejemplo verificable (muestreo uniforme sobre un cuadrado).}
Para $d=2$, $U=(-1,1)^2$ y $\Gamma=\partial U$, tomemos las
probabilidades uniformes $d\mu_r=\tfrac14\,dx$ y
$d\mu_b=\tfrac18\,d\mathcal H^1$. Si $0<\varepsilon\le1$,
particionemos cada lado del cuadrado en
$n=\lceil2/\varepsilon\rceil$ segmentos iguales y el interior en
$n^2$ cuadrados de lado $\ell=2/n$. Se cumple
$\tfrac23\varepsilon\le\ell\le\varepsilon$, de modo que el
diámetro de las celdas interiores es a lo sumo $\sqrt2\varepsilon$
y cada una tiene masa $\ell^2/4\ge\varepsilon^2/9$.
En el borde, las celdas tienen diámetro a lo sumo
$\varepsilon$ y masa $\ell/8\ge\varepsilon/12$.
Asimismo, para $0<r\le1$,
\[
\mu_r(B(x,r)\cap U)\le\tfrac\pi4 r^2,\qquad
\mu_b(B(y,r)\cap\Gamma)\le r.
\]
Por tanto, la Hipótesis~1 se cumple en este ejemplo con
$\varepsilon_0=1$, $c_r=1/9$, $C_r=\pi/4$,
$c_b=1/12$ y $C_b=1$. Estas constantes no son óptimas.
Se trata de una verificación de la \emph{geometría y la distribución}
de muestreo, no de la Hipótesis~2 sobre comparadores neuronales.

\medskip
\noindent\textbf{Observación (Relación con Shin et al.).}
Las cotas inferiores de masa local del lema auxiliar de no degeneración local no son una hipótesis adicional
frente a \cite[Assumption~3.1]{shin2020convergence}. Shin et al. formulan la no degeneración
mediante una masa mínima para cada celda de las particiones de escala $\varepsilon$; el lema
anterior muestra que esa formulación implica las cotas inferiores en bolas empleadas en esta
tesis. Se conserva así el marco geométrico original y solo se hace explícita una de sus
consecuencias.

La prueba de Shin et al. obtiene la convergencia uniforme mediante equicontinuidad,
Arzelà--Ascoli e identificación del límite. Aquí se utiliza la consecuencia local anterior para
presentar un paso directo de $L^2(\mu)$ a $L^\infty$: una función Hölder que es grande en un
punto permanece grande en una bola de radio controlado, y esa bola posee una masa mínima.
Este cambio afecta a la ruta de demostración, no a las hipótesis ni a la conclusión del
resultado de partida.
\medskip

\medskip
\noindent\textbf{Comentario técnico (geometría de muestreo).}
Las cotas inferiores de las celdas controlan la probabilidad de que la muestra cubra el dominio;
las cotas superiores de masa de bolas intervienen en la comparación entre la integral y el
promedio discreto. El lema auxiliar de no degeneración local se utilizará únicamente para convertir el control
$L^2(\mu)$, junto con una cota Hölder uniforme, en control uniforme.

\section{Generalización probabilística (Teorema 1)}
\label{thm:generalizacion_shin}

\noindent\textbf{Fuente, alcance y formulación.}
El argumento sigue \cite[Theorem~3.1 y Appendix~B]{shin2020convergence} y compara la pérdida
poblacional con la empírica mediante una cobertura simultánea del dominio y
la porción de frontera muestreada. La desigualdad es \emph{uniforme respecto del
candidato}: el evento aleatorio depende exclusivamente de los puntos de muestra,
no de la red que se elija después de observarlos. Se presentan por separado las
constantes, los pesos y el caso unidimensional para que el enunciado sea completo.

\medskip
\noindent\textbf{Enunciado (cota de generalización uniforme por cobertura).}
Supóngase la Hipótesis~1, $d\geq2$, y fíjense $\lambda_r,\lambda_b>0$ y
$0<\alpha\leq1$. Sean $m_r,m_b$ suficientemente grandes para que los radios
\begin{equation}\label{eq:thm1_radii}
\begin{aligned}
\varepsilon_r&:=\sqrt d\,c_r^{-1/d}m_r^{-1/(2d)},\\
\varepsilon_b&:=\sqrt d\,c_b^{-1/(d-1)}m_b^{-1/(2(d-1))}
\end{aligned}
\end{equation}
pertenezcan ambos a $(0,\varepsilon_0]$. Esta exigencia sobre los
\emph{radios finales}, y no únicamente sobre las escalas auxiliares
$\eta_r=\varepsilon_r/\sqrt d$ y $\eta_b=\varepsilon_b/\sqrt d$,
garantiza que las cotas superiores de masa de bolas de la Hipótesis~1
puedan utilizarse en el lema determinista.

Supóngase que $f\in C^{0,\alpha}(\overline U)$ y
$g\in C^{0,\alpha}(\Gamma)$. Consideremos candidatos $h$ tales que
$\mathcal L[h]\in C^{0,\alpha}(\overline U)$ y
$\mathcal B[h]\in C^{0,\alpha}(\Gamma)$; en particular, pueden tomarse
las redes suficientemente regulares de las clases de los teoremas
siguientes. Aquí empleamos los regularizadores \emph{exactos}
\[
\mathcal R_r(h):=[\mathcal L[h]]_{\alpha;U}^{2},\qquad
\mathcal R_b(h):=[\mathcal B[h]]_{\alpha;\Gamma}^{2}.
\]
También se admiten regularizadores mayores que estas seminormas, siempre
que se reemplace coherentemente su definición en la pérdida empírica.

Definimos primero el factor de cobertura
\begin{equation}\label{eq:thm1_Cm}
C_m:=3\max\left\{
\frac{C_r}{c_r}(\sqrt d)^d m_r^{1/2},\;
\frac{C_b}{c_b}(\sqrt d)^{d-1}m_b^{1/2}\right\}>0
\end{equation}
y los pesos residuales de regularización
\begin{equation}\label{eq:thm1_hatlambda}
\begin{aligned}
\hat\lambda^{\mathcal R}_{r,m}
&:=\frac{3\lambda_r\varepsilon_r^{2\alpha}}{C_m}
=\frac{3\lambda_r(\sqrt d)^{2\alpha}c_r^{-2\alpha/d}}{C_m}
 m_r^{-\alpha/d},\\
\hat\lambda^{\mathcal R}_{b,m}
&:=\frac{3\lambda_b\varepsilon_b^{2\alpha}}{C_m}
=\frac{3\lambda_b(\sqrt d)^{2\alpha}c_b^{-2\alpha/(d-1)}}{C_m}
 m_b^{-\alpha/(d-1)}.
\end{aligned}
\end{equation}
Para dichos pesos, la pérdida regularizada es
\begin{equation}\label{eq:loss_holder}
\begin{aligned}
\operatorname{Loss}_m(h;\lambda,\hat\lambda_m^{\mathcal R})
:=\;&\frac{\lambda_r}{m_r}\sum_{i=1}^{m_r}
\|\mathcal L[h](x_r^i)-f(x_r^i)\|^2
+\frac{\lambda_b}{m_b}\sum_{j=1}^{m_b}
\|\mathcal B[h](x_b^j)-g(x_b^j)\|^2\\
&+\hat\lambda_{r,m}^{\mathcal R}\mathcal R_r(h)
+\hat\lambda_{b,m}^{\mathcal R}\mathcal R_b(h).
\end{aligned}
\end{equation}
Sea, finalmente, la constante \emph{no negativa}
\begin{equation}\label{eq:thm1_Cprime}
C':=3(\sqrt d)^{2\alpha}
\max\left\{
\lambda_r c_r^{-2\alpha/d}[f]_{\alpha;U}^{2},\;
\lambda_b c_b^{-2\alpha/(d-1)}[g]_{\alpha;\Gamma}^{2}
\right\}.
\end{equation}

Si las muestras interior y de frontera son independientes, existe un evento
$\Omega_m$, determinado solo por ellas, con
\begin{equation}\label{eq:thm1_prob}
\mathbb P(\Omega_m)\geq
\left(1-\sqrt{m_r}\left(1-\frac1{\sqrt{m_r}}\right)^{m_r}\right)
\left(1-\sqrt{m_b}\left(1-\frac1{\sqrt{m_b}}\right)^{m_b}\right),
\end{equation}
y \emph{para todo candidato admisible $h$, simultáneamente} en
$\Omega_m$ se cumple
\begin{equation}\label{eq:thm1_main}
\boxed{
\operatorname{Loss}^{\mathrm{PINN}}(h;\lambda)
\leq C_m\operatorname{Loss}_m(h;\lambda,\hat\lambda_m^{\mathcal R})
+C'\left(m_r^{-\alpha/d}+m_b^{-\alpha/(d-1)}\right).}
\end{equation}
En particular, el enunciado permite sustituir $h$ por cualquier
minimizador empírico mensurable, aunque dependa de la muestra.
No se afirma que los radios de cobertura sean óptimos ni que el evento de
cobertura imponga por sí solo la regularidad Hölder.

\medskip
\noindent\textbf{Caso unidimensional ($d=1$, frontera completamente observada).}
Sea $U=(a,b)$ y supóngase que los datos de Dirichlet están disponibles en
ambos extremos. Definimos la contribución de frontera determinista
\[
E_{\partial U}(h):=\frac{\lambda_b}{2}
\left(|\mathcal B[h](a)-g(a)|^2+|\mathcal B[h](b)-g(b)|^2\right).
\]
La pérdida empírica no regularizada y la poblacional son, respectivamente,
\[
\begin{aligned}
\operatorname{Loss}^{\mathrm{PINN}}_{m_r,1}(h)
&:=\frac{\lambda_r}{m_r}\sum_{i=1}^{m_r}
|\mathcal L[h](x_i)-f(x_i)|^2+E_{\partial U}(h),\\
\operatorname{Loss}^{\mathrm{PINN}}_{1}(h)
&:=\lambda_r\int_U|\mathcal L[h]-f|^2\,d\mu_r+E_{\partial U}(h).
\end{aligned}
\]
Tomamos $m_r$ suficientemente grande para que
$\varepsilon_r:=c_r^{-1}m_r^{-1/2}\leq\varepsilon_0$.
El lema probabilístico implica cobertura con probabilidad al menos
$1-\sqrt{m_r}(1-m_r^{-1/2})^{m_r}$ y la parte interior del lema
determinista proporciona
\[
\begin{aligned}
\operatorname{Loss}^{\mathrm{PINN}}_{1}(h)
\leq{}&3\lambda_r C_r m_r\varepsilon_r
\frac1{m_r}\sum_{i=1}^{m_r}|\mathcal L[h](x_i)-f(x_i)|^2
+E_{\partial U}(h)\\
&+3\lambda_r\varepsilon_r^{2\alpha}
\bigl([\mathcal L[h]]_{\alpha;U}^2+[f]_{\alpha;U}^2\bigr).
\end{aligned}
\]
Con
\[
C_{m_r}^{(1)}:=\max\{1,3C_r m_r\varepsilon_r\},\qquad
\widehat\lambda_{r,m_r}^{\mathcal R}
:=\frac{3\lambda_r\varepsilon_r^{2\alpha}}{C_{m_r}^{(1)}},
\]
se obtiene, \emph{simultáneamente para todo candidato admisible},
\begin{equation}\label{eq:thm1_1d}
\begin{split}
\operatorname{Loss}^{\mathrm{PINN}}_{1}(h)
\leq{}&C_{m_r}^{(1)}\Bigl(
\operatorname{Loss}^{\mathrm{PINN}}_{m_r,1}(h)
+\widehat\lambda_{r,m_r}^{\mathcal R}
[\mathcal L[h]]_{\alpha;U}^{2}\Bigr)\\
&+3\lambda_r c_r^{-2\alpha}[f]_{\alpha;U}^2m_r^{-\alpha}.
\end{split}
\end{equation}
Asimismo,
$C_{m_r}^{(1)}=O(m_r^{1/2})$ y
$C_{m_r}^{(1)}\widehat\lambda_{r,m_r}^{\mathcal R}
=3\lambda_r c_r^{-2\alpha}m_r^{-\alpha}$.
No se regulariza la frontera en este caso, puesto que sus dos valores se
conocen exactamente; esta es la excepción unidimensional empleada por
\cite[Theorem~3.1]{shin2020convergence}.

\medskip
\noindent\textbf{Comentario técnico (regularización Hölder).}
El paso de la pérdida empírica a la pérdida teórica requiere controlar la variación del residuo entre puntos de muestreo; de lo contrario,
una pérdida pequeña en promedio no excluye picos localizados. Por eso se introduce regularización en seminormas Hölder.
Para la desigualdad de cobertura puede trabajarse tanto con seminormas de
$\mathcal L[h],\mathcal B[h]$ como con seminormas de los residuos
$r_h,b_h$, ajustando las constantes. No son penalizaciones idénticas:
pueden dar lugar a minimizadores distintos. En la formulación adoptada aquí
se regularizan los operadores y la regularidad de los datos $f,g$ aparece
separadamente en $C'$ de \eqref{eq:thm1_main}.


\medskip
\noindent\textbf{Demostración.}
Para hacer transparente la cadena lógica, separamos el argumento en dos piezas.
Primero, bajo un \emph{evento de cobertura} (la muestra es suficientemente densa en $U$ y en $\Gamma$),
demostramos una desigualdad determinista que compara la pérdida continua con la pérdida empírica y con los
términos Hölder. Segundo, probamos que dicho evento ocurre con probabilidad alta cuando $m_r,m_b$ crecen,
usando una estimación estándar de ocupación de celdas.


\medskip
\noindent\textbf{Lema de control determinista bajo un evento de cobertura (adaptado de \cite[Lemma~B.1]{shin2020convergence}).}
Sean $0<\varepsilon_r,\varepsilon_b\leq\varepsilon_0$ y suponga que las muestras satisfacen el siguiente \emph{evento de cobertura}:
\begin{itemize}
\item para todo $x\in U$ existe $x'\in T_r$ tal que $\|x-x'\|\le \varepsilon_r$;
\item para todo $y\in\Gamma$ existe $y'\in T_b$ tal que $\|y-y'\|\le \varepsilon_b$.
\end{itemize}
Entonces, para todo $h$ para el cual $\mathcal L[h]$ y $\mathcal B[h]$ poseen las seminormas H\"older finitas indicadas en el enunciado se cumple
\begin{equation}\label{eq:lemB1}
\begin{aligned}
\text{\rm Loss}^{\text{\rm PINN}}(h;\lambda)
\le\; C_m\Big[
\text{\rm Loss}^{\text{\rm PINN}}_m(h;\lambda)
+\frac{3\lambda_r\varepsilon_r^{2\alpha}}{C_m}\,[\mathcal L[h]]_{\alpha;U}^2
+\frac{3\lambda_b\varepsilon_b^{2\alpha}}{C_m}\,[\mathcal B[h]]_{\alpha;\Gamma}^2
\Big]
\\
\quad + 3\lambda_r\varepsilon_r^{2\alpha}[f]_{\alpha;U}^2
+3\lambda_b\varepsilon_b^{2\alpha}[g]_{\alpha;\Gamma}^2,
\end{aligned}
\end{equation}
donde
\[
C_m \;=\; 3\max\{\,C_r\,m_r\,\varepsilon_r^d,\;C_b\,m_b\,\varepsilon_b^{d-1}\,\}.
\]

\noindent\emph{Comentario.}
El valor de $C_m$ crece cuando la cubierta inducida por la muestra se ``degrada'' (celdas grandes);
por eso aparece multiplicando a la pérdida empírica. En el cierre del Teorema~1 se escoge
$\varepsilon_r,\varepsilon_b$ como funciones de $m_r,m_b$ de modo que $C_m$ sea del orden $m^{1/2}$.

\medskip


\noindent\emph{Prueba.}
\emph{Nota de fuente.}
La demostración sigue el argumento de \cite[Lemma~B.1]{shin2020convergence}, adaptado a nuestra notación y al caso
$\Gamma\subseteq\partial U$. El ingrediente esencial es determinista: si las muestras forman una $\varepsilon$-red en $U$ y en $\Gamma$,
entonces los residuos en puntos arbitrarios se controlan por los residuos en los puntos muestreados más cercanos, con un error de orden Hölder.

\smallskip
\textit{Paso 1: selección de vecinos cercanos y particiones inducidas.}
Bajo la hipótesis de cobertura, para cada $x\in U$ escogemos un punto $\pi_r(x)\in T_r$ tal que $\|x-\pi_r(x)\|\le \varepsilon_r$.
De modo análogo, para cada $y\in\Gamma$ escogemos $\pi_b(y)\in T_b$ con $\|y-\pi_b(y)\|\le \varepsilon_b$.
Una elección canónica es el vecino más cercano con una regla determinista de desempate, lo que produce una partición tipo Voronoi.

Definimos las celdas
\[
A_{r,i}:=\{x\in U:\ \pi_r(x)=x_r^i\},\qquad
A_{b,j}:=\{y\in\Gamma:\ \pi_b(y)=x_b^j\},
\]
y los pesos
\[
\gamma_{r,i}:=\mu_r(A_{r,i}),\qquad \gamma_{b,j}:=\mu_b(A_{b,j}).
\]
Entonces $\{A_{r,i}\}_{i=1}^{m_r}$ es una partición medible de $U$ y $\{A_{b,j}\}_{j=1}^{m_b}$ es una partición medible de $\Gamma$.

\smallskip
\noindent\textit{Paso 2: estimación puntual a partir de Hölder.}
Para $x\in U$ escribimos el residuo interior $r_h(x):=\mathcal L[h](x)-f(x)$ y descomponemos
\[
r_h(x)=\big(\mathcal L[h](x)-\mathcal L[h](\pi_r(x))\big)+\big(\mathcal L[h](\pi_r(x))-f(\pi_r(x))\big)+\big(f(\pi_r(x))-f(x)\big).
\]
Por Cauchy--Schwarz en el producto de tres copias del espacio euclídeo se obtiene
\begin{equation}\label{eq:abc_ineq}
\|a+b+c\|^2\le 3\big(\|a\|^2+\|b\|^2+\|c\|^2\big)\qquad(a,b,c\in\mathbb{R}^q).
\end{equation}
Aplicando \eqref{eq:abc_ineq} y usando las cotas Hölder,
\[
\|\mathcal L[h](x)-\mathcal L[h](\pi_r(x))\|\le [\mathcal L[h]]_{\alpha;U}\,\|x-\pi_r(x)\|^\alpha\le [\mathcal L[h]]_{\alpha;U}\,\varepsilon_r^\alpha,
\]
\[
\|f(x)-f(\pi_r(x))\|\le [f]_{\alpha;U}\,\|x-\pi_r(x)\|^\alpha\le [f]_{\alpha;U}\,\varepsilon_r^\alpha,
\]
obtenemos para todo $x\in U$:
\begin{equation}\label{eq:b1_pointwise_int}
\|r_h(x)\|^2
\le 3\|\mathcal L[h](\pi_r(x))-f(\pi_r(x))\|^2
+3\varepsilon_r^{2\alpha}[\mathcal L[h]]_{\alpha;U}^2
+3\varepsilon_r^{2\alpha}[f]_{\alpha;U}^2.
\end{equation}
De forma completamente análoga, para $y\in\Gamma$ y $b_h(y):=\mathcal B[h](y)-g(y)$,
\begin{equation}\label{eq:b1_pointwise_bd}
\|b_h(y)\|^2
\le 3\|\mathcal B[h](\pi_b(y))-g(\pi_b(y))\|^2
+3\varepsilon_b^{2\alpha}[\mathcal B[h]]_{\alpha;\Gamma}^2
+3\varepsilon_b^{2\alpha}[g]_{\alpha;\Gamma}^2.
\end{equation}

\smallskip
\noindent\textit{Paso 3: integración y reescritura como suma ponderada.}
Integramos \eqref{eq:b1_pointwise_int} respecto a $\mu_r$ y multiplicamos por $\lambda_r$:
\begin{align*}
\lambda_r\int_U \|r_h(x)\|^2\,d\mu_r(x)
&\le 3\lambda_r\int_U \|\mathcal L[h](\pi_r(x))-f(\pi_r(x))\|^2\,d\mu_r(x) \\
&\quad +3\lambda_r\varepsilon_r^{2\alpha}[\mathcal L[h]]_{\alpha;U}^2
 +3\lambda_r\varepsilon_r^{2\alpha}[f]_{\alpha;U}^2.
\end{align*}
Como $\pi_r$ es constante en cada celda $A_{r,i}$,
\[
\int_U \|\mathcal L[h](\pi_r(x))-f(\pi_r(x))\|^2\,d\mu_r(x)
=\sum_{i=1}^{m_r}\gamma_{r,i}\,\|\mathcal L[h](x_r^i)-f(x_r^i)\|^2.
\]
Análogamente, integrando \eqref{eq:b1_pointwise_bd} en $\Gamma$ respecto a $\mu_b$,
\[
\lambda_b\int_\Gamma \|b_h(y)\|^2\,d\mu_b(y)
\le 3\lambda_b\sum_{j=1}^{m_b}\gamma_{b,j}\,\|\mathcal B[h](x_b^j)-g(x_b^j)\|^2
+3\lambda_b\varepsilon_b^{2\alpha}[\mathcal B[h]]_{\alpha;\Gamma}^2
+3\lambda_b\varepsilon_b^{2\alpha}[g]_{\alpha;\Gamma}^2.
\]

\smallskip
\noindent\textit{Paso 4: cota uniforme de pesos y paso a promedios.}
Si $x\in A_{r,i}$ entonces $\|x-x_r^i\|=\|x-\pi_r(x)\|\le \varepsilon_r$, de modo que $A_{r,i}\subset B(x_r^i,\varepsilon_r)$. Por tanto
\[
\gamma_{r,i}\le \mu_r(B(x_r^i,\varepsilon_r))\le P_r^*:=\sup_{z\in U}\mu_r(B(z,\varepsilon_r)).
\]
Del mismo modo, $\gamma_{b,j}\le P_b^*:=\sup_{z\in\Gamma}\mu_b(B(z,\varepsilon_b))$.
La Hipótesis~1 provee las cotas
\[
P_r^*\le C_r\varepsilon_r^d,\qquad P_b^*\le C_b\varepsilon_b^{d-1}.
\]
En consecuencia,
\[
\sum_{i=1}^{m_r}\gamma_{r,i}\,\|\mathcal L[h](x_r^i)-f(x_r^i)\|^2
\le C_r\varepsilon_r^d\,m_r\cdot \frac{1}{m_r}\sum_{i=1}^{m_r}\|\mathcal L[h](x_r^i)-f(x_r^i)\|^2,
\]
y análogamente para el término de borde.
Sumando las contribuciones interior y de frontera y recordando la definición de $\text{\rm Loss}^{\text{\rm PINN}}_m(h;\lambda)$, se obtiene
\begin{align*}
\text{\rm Loss}^{\text{\rm PINN}}(h;\lambda)
&\le 3C_r m_r\varepsilon_r^d\,\frac{\lambda_r}{m_r}\sum_{i=1}^{m_r}\|\mathcal L[h](x_r^i)-f(x_r^i)\|^2 \\
&\quad +3C_b m_b\varepsilon_b^{d-1}\,\frac{\lambda_b}{m_b}\sum_{j=1}^{m_b}\|\mathcal B[h](x_b^j)-g(x_b^j)\|^2 \\
&\quad +\mathcal{R}_\alpha.
\end{align*}
donde $\mathcal{R}_\alpha$ es exactamente la suma de los cuatro términos Hölder (dos dependientes de $h$ y dos de los datos $f,g$) que aparecen en \eqref{eq:lemB1}.
Tomando $C_m=3\max\{C_r m_r\varepsilon_r^d,\;C_b m_b\varepsilon_b^{d-1}\}$ y reescribiendo $\mathcal{R}_\alpha$ de la forma indicada en \eqref{eq:lemB1}, se concluye el resultado. \hfill $\square$


\medskip
\noindent\textbf{Lema de cobertura con probabilidad alta (adaptado de \cite[Lemma~B.2]{shin2020convergence}).}
Sea $X\subset\mathbb{R}^d$ un conjunto acotado y medible y sea $\mu$ una medida de probabilidad concentrada en $X$.
No se exige que $X$ sea cerrado: el argumento solo utiliza las particiones medibles indicadas a continuación.
Supongamos que existe $\varepsilon_0>0$ tal que para
cada $\varepsilon\in(0,\varepsilon_0]$ se puede elegir una partición medible finita
$\{X_k^\varepsilon\}_{k=1}^{K_\varepsilon}$ de $X$, cuyos miembros son disjuntos
salvo conjuntos de $\mu$-medida nula, verificando:
\begin{enumerate}
\item[(i)] $X=\bigcup_{k=1}^{K_\varepsilon}X_k^\varepsilon$,
$\mu(X_k^\varepsilon\cap X_\ell^\varepsilon)=0$ para $k\neq\ell$, y
$\operatorname{diam}(X_k^\varepsilon)\le \sqrt d\,\varepsilon$ para todo $k$;
\item[(ii)] existen constantes $c>0$ y $s>0$ tales que $\mu(X_k^\varepsilon)\ge c\,\varepsilon^s$ para todo $k$.
\end{enumerate}
Sean $x_1,\dots,x_n$ muestras iid con ley $\mu$ y definamos el evento de cobertura
\[
A_{n,\varepsilon}:=\Big\{\forall k\in\{1,\dots,K_\varepsilon\}\ \exists i\in\{1,\dots,n\}\text{ tal que }x_i\in X_k^\varepsilon\Big\}.
\]
Entonces
\begin{equation}\label{eq:cover_prob_general}
\mathbb{P}(A_{n,\varepsilon})\ \ge\ 1-K_\varepsilon(1-c\varepsilon^s)^n\ \ge\ 1-K_\varepsilon e^{-n c\varepsilon^s}.
\end{equation}
Además, en el evento $A_{n,\varepsilon}$ se cumple que para todo $x\in X$ existe $x_j$ con
\[
\|x-x_j\|\le \sqrt d\,\varepsilon.
\]
En particular, para $n$ suficientemente grande, de modo que
$\varepsilon=c^{-1/s}n^{-1/(2s)}\leq\varepsilon_0$, se obtiene el radio
\[
\sup_{x\in X}\min_{1\le j\le n}\|x-x_j\|\le \sqrt d\,c^{-1/s}\,n^{-1/(2s)}
\]
con probabilidad al menos $1-\sqrt{n}(1-1/\sqrt n)^n$.

\noindent\emph{Prueba.}
\emph{Nota de fuente.}
La demostración sigue el argumento de cobertura de \cite[Lemma~B.2]{shin2020convergence}. La idea es estimar, vía unión de eventos,
la probabilidad de que alguna celda de una partición finita de $X$ no reciba ninguna muestra; de allí se deduce una cota para el radio de la
$\varepsilon$-red inducida por puntos i.i.d.

\smallskip
\textit{Paso 1: cota para $K_\varepsilon$.}
Por (i)--(ii),
\[
1=\mu(X)=\sum_{k=1}^{K_\varepsilon}\mu(X_k^\varepsilon)\ \ge\ K_\varepsilon\,c\,\varepsilon^s,
\]
de donde
\begin{equation}\label{eq:Kb2_bound}
K_\varepsilon\le (c\varepsilon^s)^{-1}.
\end{equation}

\smallskip
\noindent\textit{Paso 2: estimación del evento de fallo.}
Sea $E_k:=\{x_1,\dots,x_n\notin X_k^\varepsilon\}$ el evento ``la celda $k$ queda vacía''.
Si $p_k:=\mu(X_k^\varepsilon)$, por independencia
\[
\mathbb{P}(E_k)=(1-p_k)^n\le (1-c\varepsilon^s)^n\le e^{-n c\varepsilon^s}.
\]
Como $A_{n,\varepsilon}^c=\bigcup_{k=1}^{K_\varepsilon}E_k$, la desigualdad de Boole (unión) da
\[
\mathbb{P}(A_{n,\varepsilon}^c)\le \sum_{k=1}^{K_\varepsilon}\mathbb{P}(E_k)\le K_\varepsilon(1-c\varepsilon^s)^n,
\]
lo cual prueba \eqref{eq:cover_prob_general}.

\smallskip
\noindent\textit{Paso 3: del evento $A_{n,\varepsilon}$ a una cota geométrica.}
Sea $x\in X$. Elija $k$ tal que $x\in X_k^\varepsilon$. En el evento $A_{n,\varepsilon}$ existe $j$ con $x_j\in X_k^\varepsilon$.
Entonces, por (i),
\[
\|x-x_j\|\le \operatorname{diam}(X_k^\varepsilon)\le \sqrt d\,\varepsilon.
\]

\smallskip
\noindent\textit{Paso 4: elección explícita de $\varepsilon$.}
Para $n$ suficientemente grande se tiene $\varepsilon=c^{-1/s}n^{-1/(2s)}\leq\varepsilon_0$.
Con esta elección $c\varepsilon^s=n^{-1/2}$ y por \eqref{eq:Kb2_bound} se tiene $K_\varepsilon\le \sqrt n$.
Sustituyendo en \eqref{eq:cover_prob_general} obtenemos
\[
\mathbb{P}(A_{n,\varepsilon}^c)\le K_\varepsilon(1-n^{-1/2})^n\le \sqrt n\,(1-1/\sqrt n)^n,
\]
y por el Paso 3 resulta la cota del radio. \hfill $\square$


\medskip
\noindent\textbf{Cierre de la demostración del Teorema 1.}

El lema de cobertura determinista proporciona una cota \emph{determinista} entre la pérdida poblacional y la pérdida empírica
cuando los conjuntos de muestreo $T_r$ y $T_b$ son densos en $U$ y en $\Gamma$, respectivamente,
con radios $\varepsilon_r$ y $\varepsilon_b$.
La parte probabilística del Teorema~1 consiste, por tanto, en justificar que dicha densidad ocurre con alta probabilidad.

\medskip
\noindent\emph{Comentario sobre el borde.}
En \cite{shin2020convergence} el conjunto de muestreo de borde se toma sobre un subconjunto cerrado
$\Gamma\subseteq\partial U$. En esta tesis seguimos la misma convención: $\Gamma$ no se identifica con $\partial U$ en general; por ello, los términos de borde se integran y se muestrean
sobre $\Gamma$.

\smallskip
\noindent\textit{(i) Evento de cobertura en el interior.}
Aplicamos el lema probabilístico de cobertura con $X=U$, $s=d$ y $\mu=\mu_r$.
Aunque $U$ es abierto, es un conjunto acotado y medible y las particiones de la Hipótesis~1 cumplen todas las hipótesis del lema; la compacidad de $U$ no se necesita para esta estimación. Aplicamos el lema probabilístico de cobertura a la escala auxiliar $\eta_r=c_r^{-1/d}m_r^{-1/(2d)}\leq\varepsilon_0$.
El radio de cobertura resultante es
\[
\varepsilon_r=\sqrt d\,c_r^{-1/d}\,m_r^{-1/(2d)},
\]
se tiene que $T_r=\{x_r^i\}_{i=1}^{m_r}$ es $\varepsilon_r$--denso en $U$ con probabilidad al menos
$1-\sqrt{m_r}(1-1/\sqrt{m_r})^{m_r}$.

\smallskip
\noindent\textit{(ii) Evento de cobertura en el borde.}
De manera análoga, aplicamos el lema probabilístico de cobertura con $X=\Gamma$, $s=d-1$ y $\mu=\mu_b$ (medida soportada en $\Gamma$), usando la partición de borde de la Hipótesis~1 y la escala auxiliar
$\eta_b=c_b^{-1/(d-1)}m_b^{-1/(2(d-1))}\leq\varepsilon_0$.
El radio resultante es
\[
\varepsilon_b=\sqrt d\,c_b^{-1/(d-1)}\,m_b^{-1/(2(d-1))},
\]
se obtiene que $T_b=\{x_b^j\}_{j=1}^{m_b}$ es $\varepsilon_b$--denso en $\Gamma$ con probabilidad al menos
$1-\sqrt{m_b}(1-1/\sqrt{m_b})^{m_b}$.

\smallskip
\noindent\textit{(iii) Sustitución en la cota determinista.}
En la intersección de los dos eventos de cobertura anteriores, el lema de cobertura determinista da la desigualdad \eqref{eq:lemB1}.
Multiplicando los términos dentro del corchete por $C_m$ e introduciendo los pesos
$\hat{\lambda}^{\mathcal{R}}_{r,m}:=\frac{3\lambda_r}{C_m}\varepsilon_r^{2\alpha}$ y
$\hat{\lambda}^{\mathcal{R}}_{b,m}:=\frac{3\lambda_b}{C_m}\varepsilon_b^{2\alpha}$, la expresión entre corchetes coincide con
$\text{Loss}_m(h;\lambda,\hat\lambda_m^{\mathcal{R}})$ (ver \eqref{eq:loss_holder}). Además, los términos dependientes de los datos
$f$ y $g$ son exactamente
$3\lambda_r[f]_{\alpha;U}^{2}\varepsilon_r^{2\alpha}
+3\lambda_b[g]_{\alpha;\Gamma}^{2}\varepsilon_b^{2\alpha}$.
Al sustituir los radios de (i)--(ii), este sumando es menor o igual a
$C'\big(m_r^{-\alpha/d}+m_b^{-\alpha/(d-1)}\big)$ con el valor de
$C'$ fijado en el enunciado del Teorema~1. De este modo las constantes
geométricas no se contabilizan dos veces y se obtiene
\eqref{eq:thm1_main}.

\hfill $\square$

\medskip
\noindent\textbf{Uniformidad del evento.}\par
Las particiones, las ocupaciones de celdas y los radios dependen de las
muestras y de las constantes de la Hipótesis~1; no dependen de $h$.
El lema determinista, una vez fijada la cobertura, se demuestra para un
$h$ arbitrario sin imponerle restricciones adicionales de probabilidad.
Por lo tanto, la misma realización de $\Omega_m$ permite aplicar
\eqref{eq:thm1_main} a todos los $h$ admisibles, incluso a aquellos
elegidos a partir de los datos de entrenamiento.

\medskip
\noindent\textbf{Pesos de regularización para el resultado de convergencia.}
Sea $\lambda_m^{\mathcal R}=(\lambda_{r,m}^{\mathcal R},\lambda_{b,m}^{\mathcal R})$ un vector de pesos que satisface, componente a componente,
\begin{equation}\label{eq:shin_weight_condition}
\lambda_m^{\mathcal R}\geq \hat\lambda_m^{\mathcal R},
\qquad
\|\lambda_m^{\mathcal R}\|_\infty
=O\!\left(\|\hat\lambda_m^{\mathcal R}\|_\infty\right),
\end{equation}
donde $\hat\lambda_m^{\mathcal R}$ es el vector definido en el Teorema~1. En particular,
$\lambda_m^{\mathcal R}\to0$ cuando $m_r,m_b\to\infty$.
Para $d\geq2$ definimos
\begin{equation}\label{eq:shin_holder_loss}
\operatorname{Loss}_m(h;\lambda,\lambda_m^{\mathcal R})
=
\operatorname{Loss}^{\mathrm{PINN}}_m(h;\lambda)
+
\lambda_{r,m}^{\mathcal R}[\mathcal L[h]]_{\alpha;U}^{2}
+
\lambda_{b,m}^{\mathcal R}[\mathcal B[h]]_{\alpha;\Gamma}^{2}.
\end{equation}
Para el Teorema~2 fijaremos una elección concreta dentro de la familia permitida por
\eqref{eq:shin_weight_condition}. Para $d\geq2$, sea
\begin{equation}\label{eq:shin_equal_weights}
a_m:=\max\{\hat\lambda_{r,m}^{\mathcal R},\hat\lambda_{b,m}^{\mathcal R}\},
\qquad
\lambda_{r,m}^{\mathcal R}=\lambda_{b,m}^{\mathcal R}=a_m.
\end{equation}
Estos pesos cumplen $\lambda_m^{\mathcal R}\geq\hat\lambda_m^{\mathcal R}$ componente a componente y
$\|\lambda_m^{\mathcal R}\|_\infty=\|\hat\lambda_m^{\mathcal R}\|_\infty$; además, su cociente
máximo/mínimo es igual a uno. Para $d=1$ los puntos de frontera se consideran disponibles,
se omite el regularizador de borde y se toma
$\lambda_{r,m}^{\mathcal R}=\hat\lambda_{r,m}^{\mathcal R}>0$.

\medskip
\noindent\textbf{Hipótesis 2 (clases de redes; adaptación de \cite[Assumption~3.2]{shin2020convergence}).}
Sea $k$ el mayor orden de derivación que aparece en la ecuación y sea $0<\alpha\leq1$.
Supóngase que $f\in C^{0,\alpha}(\overline U)$ y $g\in C^{0,\alpha}(\Gamma)$. Para cada tamaño de muestra $m=(m_r,m_b)$, sea $\mathcal H_m$ una clase de redes neuronales tal que:
\begin{enumerate}
\item $\mathcal H_m\subset C^{k,\alpha}(\overline U)$ y, para todo $h\in\mathcal H_m$,
\[
\mathcal L[h]\in C^{0,\alpha}(\overline U),
\qquad
\mathcal B[h]\in C^{0,\alpha}(\Gamma).
\]

\item Existe $u_m^*\in\mathcal H_m$ que interpola exactamente los datos de entrenamiento, en el sentido de que
\begin{equation}\label{eq:shin_zero_empirical}
\operatorname{Loss}^{\mathrm{PINN}}_m(u_m^*;\lambda)=0.
\end{equation}

\item Los operadores aplicados a esos comparadores poseen seminormas Hölder uniformemente acotadas:
\begin{equation}\label{eq:shin_comparator_regular}
\sup_m[\mathcal L[u_m^*]]_{\alpha;U}<\infty,
\qquad
\sup_m[\mathcal B[u_m^*]]_{\alpha;\Gamma}<\infty.
\end{equation}
\end{enumerate}

\medskip
\noindent\textbf{Observación (alcance de la Hipótesis 2).}
Esta es la hipótesis utilizada por Shin et al. en su resultado principal.
Los comparadores pueden depender de la muestra, pero la cota de sus
seminormas utilizada en la demostración debe ser uniforme a lo largo
de la sucesión considerada (con constante determinista o, si se adopta
una formulación casi segura, con constante finita en el evento correspondiente).
No se deduce de una aproximación universal sin una construcción adicional. No exige que $u_m^*$ coincida con la solución exacta $u^*$ ni que converja a ella en $C^2$; exige que la clase $\mathcal H_m$ contenga un interpolante exacto de los datos finitos con regularidad Hölder controlada. Si $u^*\in\mathcal H_m$ para todo $m$, las condiciones anteriores se satisfacen tomando $u_m^*=u^*$. El artículo señala además la posibilidad de relajar cuantitativamente la condición \eqref{eq:shin_zero_empirical}. En la Sección~\ref{sec:comparadores_aproximados} desarrollamos y demostramos una formulación explícita de esa relajación, incluida la suboptimalidad del entrenamiento; no la atribuimos como un resultado de originalidad establecida.

\section{Convergencia de la pérdida y de los residuos (Teorema 2)}

\noindent\textbf{Enunciado (pérdida poblacional y convergencia uniforme de los operadores).}
Supónganse las Hipótesis~1 y~2 y fíjense pesos residuales
$\lambda_r,\lambda_b>0$, independientes de $m$. Para $d\geq2$, sean $m_r,m_b$ los números de muestras
independientes e i.i.d. en el interior y el borde, respectivamente, con
\[
m_r=O\!\left(m_b^{d/(d-1)}\right).
\]
Elijamos los pesos de regularización \emph{iguales} dados por
\eqref{eq:shin_equal_weights}, y sea
\[
h_m\in\operatorname*{arg\,min}_{h\in\mathcal H_m}
\operatorname{Loss}_m(h;\lambda,\lambda_m^{\mathcal R})
\]
un minimizador global de \eqref{eq:shin_holder_loss}.
Se presupone la existencia de una elección mensurable de minimizadores y
se interpreta el muestreo para todos los tamaños sobre un mismo espacio
de probabilidad (por ejemplo, mediante dos sucesiones infinitas de muestras
i.i.d., interior y de frontera). Las conclusiones casi seguras se consideran
a lo largo de una sucesión
$(m_r(n),m_b(n))$ de tamaños de muestra que satisface
$m_r(n),m_b(n)\geq n$ para $n$ suficientemente grande.
Entonces:
\begin{enumerate}
\item en el evento de cobertura del Teorema~1,
\begin{equation}\label{eq:shin_expected_rate}
\operatorname{Loss}^{\mathrm{PINN}}(h_m;\lambda)
=O\!\left(m_r^{-\alpha/d}\right);
\end{equation}
\item con probabilidad uno respecto del muestreo i.i.d.,
\begin{equation}\label{eq:shin_uniform_operators}
\mathcal L[h_m]\longrightarrow f
\quad\text{en }C^0(\overline U),
\qquad
\mathcal B[h_m]\longrightarrow g
\quad\text{en }C^0(\Gamma).
\end{equation}
En el mismo evento, la demostración proporciona además las cotas
\begin{equation}\label{eq:shin_sup_rates}
\begin{aligned}
\|\mathcal L[h_m]-f\|_{C^0(\overline U)}
&=O\!\left(m_r^{-\alpha^2/[d(2\alpha+d)]}\right),\\
\|\mathcal B[h_m]-g\|_{C^0(\Gamma)}
&=O\!\left(m_r^{-\alpha^2/[d(2\alpha+d-1)]}\right).
\end{aligned}
\end{equation}
Estas tasas son consecuencias de la desigualdad de pérdida y de la
regularidad Hölder uniforme; no son tasas de aproximación óptimas ni
tasas de entrenamiento de los optimizadores numéricos.
\end{enumerate}
Para $d=1$, a lo largo de tamaños $m_r(n)\geq n$ y con el único peso
$\lambda_{r,m}^{\mathcal R}=\hat\lambda_{r,m}^{\mathcal R}$, se obtienen las
conclusiones correspondientes, con tasa $O(m_r^{-\alpha})$ y con todos los datos de
frontera disponibles. El argumento unidimensional se detalla al final de la demostración.

\noindent\emph{Alcance del enunciado.}
Esta es una versión con elección explícita de pesos del resultado de
\cite[Theorem~3.2]{shin2020convergence}; no se afirma aquí el resultado para
\emph{todos} los vectores que satisfacen \eqref{eq:shin_weight_condition}.

\medskip
\noindent\textbf{Demostración.}
La primera parte reproduce el argumento del Apéndice~C de Shin et al.; la última etapa, que mejora la convergencia $L^2$ a convergencia uniforme, se presenta mediante el lema de no degeneración local deducido de las particiones.

\smallskip
\noindent\emph{Paso 1: control uniforme de las seminormas regularizadas.}
Para $d\geq2$, definamos
\[
R_r(h):=[\mathcal L[h]]_{\alpha;U}^{2},
\qquad
R_b(h):=[\mathcal B[h]]_{\alpha;\Gamma}^{2},
\]
y
\[
R^*:=\sup_m\bigl(R_r(u_m^*)+R_b(u_m^*)\bigr)<\infty,
\]
por \eqref{eq:shin_comparator_regular}. Como $h_m$ minimiza la pérdida regularizada,
$u_m^*\in\mathcal H_m$ interpola los datos y los pesos son iguales a $a_m>0$,
\begin{align*}
a_m\bigl(R_r(h_m)+R_b(h_m)\bigr)
&\leq \operatorname{Loss}_m(h_m;\lambda,\lambda_m^{\mathcal R})\\
&\leq \operatorname{Loss}_m(u_m^*;\lambda,\lambda_m^{\mathcal R})\\
&=a_m\bigl(R_r(u_m^*)+R_b(u_m^*)\bigr)\\
&\leq a_m R^*.
\end{align*}
Dividiendo por $a_m$, obtenemos directamente
\begin{equation}\label{eq:shin_uniform_regularizers}
R_r(h_m)+R_b(h_m)\leq R^*,
\end{equation}
sin presuponer la comparabilidad de dos pesos distintos. En particular,
\begin{equation}\label{eq:shin_empirical_loss_rate}
\operatorname{Loss}_m(h_m;\lambda,\lambda_m^{\mathcal R})\leq a_mR^*.
\end{equation}

\smallskip
\noindent\emph{Paso 2: tasa para la pérdida poblacional.}
Escribamos
\[
A_r:=3\lambda_r(\sqrt d)^{2\alpha}c_r^{-2\alpha/d},
\qquad
A_b:=3\lambda_b(\sqrt d)^{2\alpha}c_b^{-2\alpha/(d-1)}.
\]
Por la definición de $\hat\lambda_m^{\mathcal R}$ en el Teorema~1 y
la elección \eqref{eq:shin_equal_weights},
\begin{equation}\label{eq:shin_Cmam_rate}
C_ma_m=\max\{A_rm_r^{-\alpha/d},\,A_bm_b^{-\alpha/(d-1)}\}.
\end{equation}
De $m_r=O(m_b^{d/(d-1)})$ se deduce
$m_b^{-\alpha/(d-1)}=O(m_r^{-\alpha/d})$; en consecuencia
\[
C_ma_m=O(m_r^{-\alpha/d}).
\]
Además, $C_m\geq c\sqrt{m_r}$ para cierta constante $c>0$ independiente de $m$,
lo que, combinado con la cota anterior, da
$a_m=O(m_r^{-1/2-\alpha/d})$. Por
\eqref{eq:shin_empirical_loss_rate}, esta es también una cota superior para la
pérdida regularizada empírica (multiplicada por $R^*$).

El evento de cobertura del Teorema~1 depende exclusivamente de las muestras,
y no de la función candidata $h$. Por ello, la desigualdad probabilística
vale \emph{simultáneamente} para todas las funciones con seminormas finitas,
incluido el minimizador aleatorio $h_m$ elegido a partir de esas mismas
muestras. En dicho evento, puesto que
$\lambda_m^{\mathcal R}\geq\hat\lambda_m^{\mathcal R}$ componente a componente,
\begin{align*}
\operatorname{Loss}^{\mathrm{PINN}}(h_m;\lambda)
&\leq C_m\operatorname{Loss}_m(h_m;\lambda,\hat\lambda_m^{\mathcal R})
 +C'\bigl(m_r^{-\alpha/d}+m_b^{-\alpha/(d-1)}\bigr)\\
&\leq C_ma_mR^*
 +C'\bigl(m_r^{-\alpha/d}+m_b^{-\alpha/(d-1)}\bigr)\\
&=O(m_r^{-\alpha/d}),
\end{align*}
lo que demuestra \eqref{eq:shin_expected_rate}. Consideremos una sucesión de pares de tamaños muestrales enteros
$(m_r(n),m_b(n))$ con $m_r(n)\geq n$ y $m_b(n)\geq n$ para $n$
suficientemente grande, que satisfaga la relación anterior. Las cotas de fallo
del Teorema~1 son sumables a lo largo de esta sucesión: cada una es de orden
$\sqrt m\,(1-m^{-1/2})^m\leq\sqrt m\,e^{-\sqrt m}$,
una función decreciente y sumable sobre los enteros grandes. Por Borel--Cantelli, con probabilidad uno la estimación vale finalmente y, como
$\lambda_r,\lambda_b>0$ son constantes fijas, se sigue que
\begin{equation}\label{eq:shin_L2_convergence}
\|\mathcal L[h_m]-f\|_{L^2(U;\mu_r)}\to0,
\qquad
\|\mathcal B[h_m]-g\|_{L^2(\Gamma;\mu_b)}\to0.
\end{equation}

\smallskip
\noindent\emph{Paso 3: control cuantitativo de $L^2(\mu)$ a $C^0$.}
De \eqref{eq:shin_uniform_regularizers} y de la desigualdad triangular
para seminormas Hölder se obtiene
\begin{equation}\label{eq:uniform_holder_residual}
[\mathcal L[h_m]-f]_{\alpha;U}
\leq\sqrt{R^*}+[f]_{\alpha;U}=:K_r,
\qquad
[\mathcal B[h_m]-g]_{\alpha;\Gamma}
\leq\sqrt{R^*}+[g]_{\alpha;\Gamma}=:K_b.
\end{equation}
Estas constantes no dependen del tamaño de la muestra.

\medskip
\noindent\textit{Lema auxiliar (interpolación local Hölder--$L^2$).}
Sea $(X,\mathrm{dist},\mu)$ un espacio métrico de medida de probabilidad
con $\mu(B(x,\rho))\geq c\rho^s$ para $x\in X$ y
$0<\rho\leq\rho_0$, donde $c,\rho_0,s>0$. Si $v:X\to\mathbb R$
es acotada y $[v]_{C^{0,\alpha}(X)}\leq K$, entonces existe una
constante $C=C(c,\rho_0,s,\alpha,K)$ tal que
\begin{equation}\label{eq:holder_L2_interpolacion}
\|v\|_{L^\infty(X)}\leq
C\left(\|v\|_{L^2(X;\mu)}+
\|v\|_{L^2(X;\mu)}^{\frac{2\alpha}{2\alpha+s}}\right).
\end{equation}
\emph{Demostración del lema.}
Pongamos $M=\|v\|_\infty$. Si $M=0$ no hay nada que probar.
El caso $K=0$ corresponde a funciones constantes y da
$M=\|v\|_{L^2(\mu)}$. Si $K>0$, escogemos $x\in X$ tal que
$|v(x)|\geq3M/4$, y definimos
\[
\rho:=\min\!\left\{\rho_0/2,\left(\frac{M}{4K}\right)^{1/\alpha}\right\}.
\]
Para $z\in B(x,\rho)$, la condición Hölder implica
$|v(z)|\geq M/2$. Por la cota inferior de masa,
\begin{equation}\label{eq:holder_L2_mass}
\|v\|_{L^2(\mu)}^2\geq
\frac{cM^2}{4}\min\!\left\{(\rho_0/2)^s,
\left(\frac{M}{4K}\right)^{s/\alpha}\right\}.
\end{equation}
Si el primer término es el mínimo, obtenemos
$M\leq C_1\|v\|_{L^2(\mu)}$. Si lo es el segundo,
$M^{2+s/\alpha}\leq C_2\|v\|_{L^2(\mu)}^2$, de donde
$M\leq C_3\|v\|_{L^2(\mu)}^{2\alpha/(2\alpha+s)}$.
La combinación de ambos casos da \eqref{eq:holder_L2_interpolacion}.
\hfill$\square$

\medskip
Aplicamos el lema a $X=U$, $\mu=\mu_r$, $s=d$, $v=\mathcal L[h_m]-f$;
y a $X=\Gamma$, $\mu=\mu_b$, $s=d-1$,
$v=\mathcal B[h_m]-g$. La hipótesis sobre las masas se sigue del
lema auxiliar de no degeneración local y las cotas Hölder uniformes son
\eqref{eq:uniform_holder_residual}. Por
\eqref{eq:shin_L2_convergence}, ambas normas supremas tienden a cero.
Como los operadores y los datos poseen extensiones continuas al cierre,
esto demuestra \eqref{eq:shin_uniform_operators}.

Para obtener las tasas anunciadas, la no negatividad de los dos términos
de la pérdida poblacional y $\lambda_r,\lambda_b>0$ dan
\[
\|\mathcal L[h_m]-f\|_{L^2(\mu_r)}^2,
\ \|\mathcal B[h_m]-g\|_{L^2(\mu_b)}^2
=O(m_r^{-\alpha/d}),
\]
con constantes que pueden depender de los pesos residuales fijos.
Sustituyendo en \eqref{eq:holder_L2_interpolacion} con
$s=d$ y $s=d-1$, respectivamente, obtenemos
\eqref{eq:shin_sup_rates}. En particular, las conclusiones son casi
seguras a lo largo de la sucesión fijada, pues el argumento
probabilístico del Paso~2 ya produjo un único evento de probabilidad uno.
\hfill$\square$

\medskip
\noindent\emph{Caso $d=1$: pérdida, cota y convergencia uniforme.}
Aquí $U=(a,b)$ y los dos valores de frontera se evalúan sin muestreo. Utilizamos
exactamente las pérdidas y el peso $\widehat\lambda_{r,m_r}^{\mathcal R}$
definidos en \eqref{eq:thm1_1d}; la pérdida regularizada es
\[
\operatorname{Loss}^{\mathcal R}_{m_r,1}(h)
:=\operatorname{Loss}^{\mathrm{PINN}}_{m_r,1}(h)
 +\widehat\lambda_{r,m_r}^{\mathcal R}
  [\mathcal L[h]]_{\alpha;U}^2.
\]
Sea $R_r^*:=\sup_{m_r}[\mathcal L[u_{m_r}^*]]_{\alpha;U}^2<\infty$.
Por minimalidad e interpolación exacta,
\[
\begin{aligned}
\widehat\lambda_{r,m_r}^{\mathcal R}
 [\mathcal L[h_{m_r}]]_{\alpha;U}^2
&\leq \operatorname{Loss}^{\mathcal R}_{m_r,1}(h_{m_r})\\
&\leq \operatorname{Loss}^{\mathcal R}_{m_r,1}(u_{m_r}^*)
\leq \widehat\lambda_{r,m_r}^{\mathcal R}R_r^*.
\end{aligned}
\]
En consecuencia,
$[\mathcal L[h_{m_r}]]_{\alpha;U}^2\leq R_r^*$ y
$\operatorname{Loss}^{\mathcal R}_{m_r,1}(h_{m_r})
\leq\widehat\lambda_{r,m_r}^{\mathcal R}R_r^*$.
Aplicando \eqref{eq:thm1_1d} en el evento de cobertura obtenemos
\[
\begin{aligned}
\operatorname{Loss}^{\mathrm{PINN}}_{1}(h_{m_r})
&\leq C_{m_r}^{(1)}\operatorname{Loss}^{\mathcal R}_{m_r,1}(h_{m_r})
+3\lambda_r c_r^{-2\alpha}[f]_{\alpha;U}^{2}m_r^{-\alpha}\\
&\leq 3\lambda_r c_r^{-2\alpha}
\bigl(R_r^*+[f]_{\alpha;U}^{2}\bigr)m_r^{-\alpha}.
\end{aligned}
\]
Por tanto el término interior poblacional converge a cero en $L^2(\mu_r)$
y la contribución $E_{\partial U}(h_{m_r})$ también tiende a cero.
Como $\lambda_b>0$, cada uno de los dos defectos de frontera converge a cero.

Para completar el paso de $L^2$ a uniforme sin invocar el lema auxiliar,
que arriba se formuló para $d\geq2$, observamos que las particiones de
la Hipótesis~1 con $d=1$ implican, para todo $x\in U$ y
$0<\rho\leq\varepsilon_0$, la cota local
\[
\mu_r(B(x,\rho)\cap U)\geq \tfrac12c_r\rho.
\]
En efecto, basta escoger una celda de la partición de escala $\rho/2$
que contenga a $x$; toda esa celda tiene diámetro a lo sumo $\rho/2$
y masa al menos $c_r\rho/2$. Por otro lado,
\[
[\mathcal L[h_{m_r}]-f]_{\alpha;U}
\leq \sqrt{R_r^*}+[f]_{\alpha;U}.
\]
Si el residuo fuese mayor que un $\delta>0$ en algún punto de una
subsucesión, esta cota Hölder lo mantendría mayor que $\delta/2$ en
$B(x,\rho)\cap U$ para un $\rho>0$ independiente de $m_r$.
La cota inferior de masa daría entonces una cota positiva para su norma
$L^2(\mu_r)$, en contradicción con la convergencia anterior.
Así, $\mathcal L[h_{m_r}]\to f$ uniformemente en $\overline U$.

Finalmente, para una sucesión de tamaños $m_r(n)\geq n$, las probabilidades
de fallo $\sqrt{m_r}(1-m_r^{-1/2})^{m_r}$ son sumables. El lema de
Borel--Cantelli convierte las cotas obtenidas en conclusiones casi seguras
a lo largo de esa sucesión.

\medskip
\noindent\textbf{Observación (qué tasa se demuestra).}
La expresión $O(m_r^{-\alpha/d})$ corresponde a una pérdida residual
\emph{cuadrática}, no a la norma uniforme del residuo ni al error
$\|h_m-u^*\|$. La cota \eqref{eq:shin_sup_rates} es una consecuencia
adicional, más débil, del control Hölder uniforme y de las masas de bolas.
No se extrapola de estas estimaciones una pendiente $-1$ para errores
numéricos ni se afirma convergencia de los procedimientos Adam o L-BFGS.

\medskip
\noindent\textbf{Observación (alcance de la condición sobre pesos).}
La condición \eqref{eq:shin_weight_condition}, por sí sola, no garantiza que
$\max\{\lambda_{r,m}^{\mathcal R},\lambda_{b,m}^{\mathcal R}\}/
\min\{\lambda_{r,m}^{\mathcal R},\lambda_{b,m}^{\mathcal R}\}$ permanezca
acotado. Por ejemplo, si $d=2$ y $m_r=m_b=n$, los pesos mínimos del
Teorema~1 verifican, salvo constantes independientes de $n$,
$\hat\lambda_{r,m}^{\mathcal R}\asymp n^{-1/2-\alpha/2}$ y
$\hat\lambda_{b,m}^{\mathcal R}\asymp n^{-1/2-\alpha}$;
su cociente crece como $n^{\alpha/2}$. La elección
\eqref{eq:shin_equal_weights} evita esa dificultad. De este modo se prueba
la conclusión para una familia particular de pesos permitida por Shin, sin
atribuir a las condiciones generales una consecuencia que no se desprende de ellas.

\medskip
\par\medskip
\noindent\textbf{Observación (minimizadores exactos y entrenamiento práctico).}
El Teorema~2 adapta el marco de Shin et al. con la elección explícita de pesos \eqref{eq:shin_equal_weights} (para $d\geq2$) y supone minimizadores globales exactos. El propio artículo contempla aproximaciones a la minimización y la posibilidad de relajar la interpolación exacta. La Sección~\ref{sec:comparadores_aproximados} determina una escala suficiente para ambas tolerancias sin modificar el enunciado precedente, que sigue siendo la reconstrucción del resultado de Shin para pesos iguales.
\par\medskip
\noindent\textbf{Observación (Paso al problema de Dirichlet clásico).}
En los resultados que siguen fijamos $\Gamma=\partial U$; es decir, el dato de contorno $g$ se prescribe en toda la frontera.
Esta identificación es necesaria para la unicidad del problema de Dirichlet y para las estimaciones de estabilidad que
convierten el control del residuo en control del error $h-u^*$.
\medskip


\section{Extensión condicional: comparadores y minimizadores aproximados}
\label{sec:comparadores_aproximados}

La Hipótesis~2 impone un comparador que anula exactamente la pérdida
empírica. Esta exigencia es distinta de la mera capacidad de aproximación de
una clase neuronal. El argumento anterior permite formular una variante que
acepta pérdida residual empírica positiva y minimización subóptima, siempre
que ambas tolerancias sean controladas en relación con el peso regularizador.
La posibilidad de considerar soluciones subóptimas y comparadores de pérdida
no nula ya se discute en el marco de \cite{shin2020convergence}; el propósito
de la siguiente proposición es explicitar una condición cuantitativa
\emph{suficiente} y demostrar su compatibilidad con las estimaciones de este
capítulo. No se afirma originalidad bibliográfica de la idea general.

\medskip
\noindent\textbf{Proposición (comparadores aproximados y tolerancia de optimización).}
Supónganse las condiciones de muestreo y regularidad de la Hipótesis~1,
con $d\geq2$ y $0<\alpha\leq1$, $f\in C^{0,\alpha}(\overline U)$,
$g\in C^{0,\alpha}(\Gamma)$ y pesos residuales fijos
$\lambda_r,\lambda_b>0$; se requieren asimismo las condiciones del
apartado 1 de la Hipótesis~2 para las clases $\mathcal H_m$. Sean $m_r,m_b$ los tamaños
muestrales del Teorema~2, con
$m_r=O(m_b^{d/(d-1)})$, y elíjanse los pesos iguales
$a_m>0$ de \eqref{eq:shin_equal_weights}. Para cada par $m$, sean
$v_m,h_m\in\mathcal H_m$ elecciones mensurables, eventualmente
aleatorias, tales que
\begin{align}
\operatorname{Loss}^{\mathrm{PINN}}_m(v_m;\lambda)
 &\leq\varepsilon_m, \label{eq:aprox_comparador_loss}\\
R_r(v_m)+R_b(v_m)&\leq R^*,\label{eq:aprox_comparador_holder}\\
\operatorname{Loss}_m(h_m;\lambda,(a_m,a_m))
 &\leq \inf_{h\in\mathcal H_m}
 \operatorname{Loss}_m(h;\lambda,(a_m,a_m))+\delta_m.
 \label{eq:aprox_minimizador}
\end{align}
Aquí $\varepsilon_m,\delta_m\geq0$ son cotas deterministas,
$R^*<\infty$ no depende de $m$, y las desigualdades se suponen
válidas casi seguramente, al menos para todos los índices suficientemente
grandes de la sucesión de tamaños. Los regularizadores son
$R_r(h)=[\mathcal L[h]]_{\alpha;U}^2$ y
$R_b(h)=[\mathcal B[h]]_{\alpha;\Gamma}^2$.

Entonces, con $\kappa_m:=(\varepsilon_m+\delta_m)/a_m$, se tiene
\begin{align}
R_r(h_m)+R_b(h_m)&\leq R^*+\kappa_m,
\label{eq:aprox_regularidad_resultado}\\
\operatorname{Loss}_m(h_m;\lambda,(a_m,a_m))
 &\leq a_m(R^*+\kappa_m),
\label{eq:aprox_loss_empirica}\\
\operatorname{Loss}^{\mathrm{PINN}}(h_m;\lambda)
 &\leq C_ma_m(R^*+\kappa_m)
 +C'\bigl(m_r^{-\alpha/d}+m_b^{-\alpha/(d-1)}\bigr)
\label{eq:aprox_loss_poblacional}
\end{align}
en el evento de cobertura del Teorema~1. En particular, si
\begin{equation}\label{eq:aprox_escala_suficiente}
\boxed{\qquad\varepsilon_m+\delta_m\leq K a_m
\quad\text{para todo $m$ suficientemente grande,}\qquad}
\end{equation}
con una constante determinista $K<\infty$, entonces, con probabilidad
uno a lo largo de la sucesión muestral del Teorema~2,
\begin{equation}\label{eq:aprox_tasa_poblacional}
\operatorname{Loss}^{\mathrm{PINN}}(h_m;\lambda)
=O(m_r^{-\alpha/d}),
\end{equation}
y se mantienen las convergencias uniformes y las tasas de
\eqref{eq:shin_uniform_operators}--\eqref{eq:shin_sup_rates}.

\medskip
\noindent\textit{Demostración.}
Por no negatividad de la pérdida PINN empírica y por la
$\delta_m$-suboptimalidad,
\begin{align*}
a_m\bigl(R_r(h_m)+R_b(h_m)\bigr)
&\leq\operatorname{Loss}_m(h_m;\lambda,(a_m,a_m))\\
&\leq\operatorname{Loss}_m(v_m;\lambda,(a_m,a_m))+\delta_m\\
&=\operatorname{Loss}^{\mathrm{PINN}}_m(v_m;\lambda)
   +a_m(R_r(v_m)+R_b(v_m))+\delta_m\\
&\leq \varepsilon_m+a_mR^*+\delta_m.
\end{align*}
La división por $a_m>0$ demuestra
\eqref{eq:aprox_regularidad_resultado}; la misma comparación demuestra
\eqref{eq:aprox_loss_empirica}.

El evento de cobertura es uniforme respecto de todas las funciones
admisibles, incluidas las que dependen de los datos. Por el Teorema~1 y
$\widehat\lambda_{r,m}^{\mathcal R},
\widehat\lambda_{b,m}^{\mathcal R}\leq a_m$,
\begin{align*}
\operatorname{Loss}^{\mathrm{PINN}}(h_m;\lambda)
&\leq C_m\operatorname{Loss}_m
 (h_m;\lambda,\widehat\lambda_m^{\mathcal R})
 +C'\bigl(m_r^{-\alpha/d}+m_b^{-\alpha/(d-1)}\bigr)\\
&\leq C_m\bigl(\varepsilon_m+\delta_m+a_mR^*\bigr)
 +C'\bigl(m_r^{-\alpha/d}+m_b^{-\alpha/(d-1)}\bigr),
\end{align*}
que equivale a \eqref{eq:aprox_loss_poblacional}.
Si se satisface \eqref{eq:aprox_escala_suficiente}, las seminormas de
$\mathcal L[h_m]$ y $\mathcal B[h_m]$ están acotadas uniformemente por
$R^*+K$. Por \eqref{eq:shin_Cmam_rate},
$C_ma_m=O(m_r^{-\alpha/d})$; la relación entre $m_r$ y $m_b$
permite absorber el término de frontera. Así se obtiene
\eqref{eq:aprox_tasa_poblacional}. Las probabilidades de fallo de
cobertura son sumables, según el Paso~2 de la demostración del
Teorema~2; Borel--Cantelli proporciona un único evento de probabilidad
uno donde las cotas valen para todos los índices suficientemente grandes.
Por último, las seminormas Hölder uniformes y el lema de interpolación
\eqref{eq:holder_L2_interpolacion} permiten convertir la convergencia
$L^2(\mu_r)$ y $L^2(\mu_b)$ en convergencia uniforme, con los mismos
exponentes de \eqref{eq:shin_sup_rates}. \hfill$\square$

\medskip
\noindent\textbf{Corolario (criterio suficiente mediante aproximación fuerte).}
Sea $u^*\in C^{2,\alpha}(\overline U)$ solución exacta de un problema
elíptico lineal de segundo orden con contorno Dirichlet, y supóngase que
$\mathcal L:C^{2,\alpha}(\overline U)\to C^{0,\alpha}(\overline U)$
y $\mathcal B:C^{2,\alpha}(\overline U)\to C^{0,\alpha}(\partial U)$
son aplicaciones lineales continuas, con normas de operador
independientes de $m$ (por ejemplo, bajo regularidad global suficiente
de sus coeficientes). Si para todo $m$ existe $v_m\in\mathcal H_m$
con
\begin{equation}\label{eq:aprox_criterio_fuerte}
\|v_m-u^*\|_{C^{2,\alpha}(\overline U)}
\leq A\sqrt{a_m},
\end{equation}
para una constante $A$ independiente de $m$, entonces los
comparadores $v_m$ cumplen
\eqref{eq:aprox_comparador_loss} y
\eqref{eq:aprox_comparador_holder} con
$\varepsilon_m=O(a_m)$ y $R^*<\infty$, de modo que la proposición
precedente se aplica a cualquier sucesión de soluciones
$\delta_m$-subóptimas con $\delta_m=O(a_m)$.

\noindent\textit{Demostración.}
Como $\mathcal L[u^*]=f$ y $\mathcal B[u^*]=g$,
la continuidad de los operadores y
\eqref{eq:aprox_criterio_fuerte} dan
$\|\mathcal L[v_m]-f\|_\infty
+\|\mathcal B[v_m]-g\|_\infty=O(\sqrt{a_m})$.
Al elevar al cuadrado y evaluar en las muestras, se obtiene
$\operatorname{Loss}^{\mathrm{PINN}}_m(v_m;\lambda)=O(a_m)$,
con constantes independientes de las muestras.
Por otra parte, la desigualdad triangular y la continuidad de los
operadores en norma Hölder implican
$[\mathcal L[v_m]]_{\alpha;U}+[\mathcal B[v_m]]_{\alpha;\partial U}
\leq [f]_{\alpha;U}+[g]_{\alpha;\partial U}+O(\sqrt{a_m})$.
Dado que $a_m\to0$, estas seminormas están uniformemente acotadas.
\hfill$\square$

\medskip
\noindent\textbf{Consecuencias para los Teoremas 3 y 4.}
Si se suponen además el contorno Dirichlet completo y las condiciones
elípticas de estabilidad del Teorema~3, los aproximantes de la
proposición satisfacen su conclusión de convergencia en $C^0$ y la
tasa allí indicada. Para Poisson, si $U$ es Lipschitz y
$u^*,h_m\in H^1(U)$ tienen la misma traza en toda la frontera, la
estimación energética del Teorema~4 proporciona convergencia en
$H^1(U)$. Bajo la cota inferior de densidad interior incluida en la
Hipótesis~1, se conserva la tasa condicional
$O(m_r^{-\alpha/(2d)})$.
Estas consecuencias requieren escoger las clases $\mathcal H_m$ de modo
compatible con Dirichlet fuerte cuando este se impone: la aproximación
\eqref{eq:aprox_criterio_fuerte}, por sí sola, no impone igualdad exacta
de trazas.

\medskip
\noindent\textbf{Alcance y límites.}
La condición $\varepsilon_m\to0$ por sí sola no permite cerrar la
prueba: el cociente $(\varepsilon_m+\delta_m)/a_m$ podría divergir,
y entonces \eqref{eq:aprox_regularidad_resultado} no ofrece una cota
Hölder uniforme. La condición \eqref{eq:aprox_criterio_fuerte} es un
\emph{criterio suficiente}, no un teorema que garantice la existencia
de redes con dicha tasa de aproximación para todas las arquitecturas.
Asimismo, una secuencia de pérdidas pequeñas obtenida con Adam o
L-BFGS no demuestra que la brecha de suboptimalidad $\delta_m$ sea
$O(a_m)$, pues no se dispone de una cota global de optimalidad para
esos entrenamientos. Esta proposición no asigna tasas de convergencia a
los optimizadores del Colab.

\section{Existencia y regularidad para el problema de Dirichlet}
\label{sec:existencia_dirichlet}

(Teorema 6.13 en \cite{gilbarg2001elliptic}.)
Sea $0<\alpha<1$ y sea $U$ un dominio acotado que satisfaga la condición de esfera exterior en cada punto de
$\partial U$. Supóngase que el operador $\mathcal{L}$ (Capítulo~1) es uniformemente elíptico en $U$,
que sus coeficientes satisfacen las hipótesis de regularidad y acotación de la Hipótesis~3 (abajo) y que
\[
\tilde c(x):=\sum_{i=1}^d D_i b^i(x)+d(x)\le 0 \quad \text{en } U.
\]
Entonces, para $f\in C^{0,\alpha}(U)$ acotada y $g\in C^0(\partial U)$, el problema de Dirichlet del Capítulo~1 admite una única solución clásica
$u^*\in C^{2,\alpha}_{\mathrm{loc}}(U)\cap C^0(\overline U)$.
La notación $C^{2,\alpha}_{\mathrm{loc}}(U)$ expresa regularidad de segundo orden en todo compacto contenido en $U$;
la continuidad de $g$ no implica, por sí sola, una cota global de la norma $C^{2,\alpha}$ hasta la frontera.

\medskip
\noindent\textbf{Comentario técnico (estabilidad del operador).}
Este resultado garantiza que el objeto límite $u^*$ existe y es regular, lo que permite cerrar el
argumento de convergencia de PINNs. La condición $\tilde c\le 0$ (en combinación con elipticidad uniforme)
es la que habilita un principio del máximo y, por tanto, estimaciones de estabilidad en norma suprema,
ver estimación de estabilidad de la Sección~\ref{sec:estabilidad_dirichlet}.

\section{Condiciones de elipticidad y regularidad (Hipótesis 3)}

Consideremos $\lambda(x)$ el valor propio mínimo de $[a^{ij}(x)]$ y $\alpha\in(0,1)$.

\begin{enumerate}
\item (\textit{Elipticidad uniforme}) Existe $\lambda_0>0$ tal que $\lambda(x)\ge \lambda_0$ para todo $x\in U$
y $a^{ij}=a^{ji}$.
\item Los coeficientes $c^i$ y $d$ pertenecen a $C^{0,\alpha}(U)$ y tienen normas $L^\infty(U)$ finitas.
\item Los coeficientes $a^{ij}$ y $b^i$ pertenecen a $C^{1,\alpha}(U)$, con coeficientes y derivadas de primer orden acotados en $U$. Además, $f\in C^{0,\alpha}(U)$ es acotada y $g\in C^{0,\alpha}(\partial U)$. Estas condiciones aseguran, en particular, que $\widetilde b$ y $\widetilde c$ están acotados, además de poseer la regularidad H\"older local que requiere el problema clásico de Dirichlet.
\item $U$ satisface la condición de esfera exterior en cada punto de la frontera.
\item Existen constantes $\Lambda,\nu\ge 0$ tales que, para todo $x\in U$,
\[
\sum_{i,j}|a^{ij}(x)|^2\le \Lambda,\qquad
\lambda_0^{-2}\big(\|b(x)\|^2+\|c(x)\|^2\big)+\lambda_0^{-1}|d(x)|\le \nu^2,
\]
donde $b(x)=[b^1(x),\dots,b^d(x)]^T$ y $c(x)=[c^1(x),\dots,c^d(x)]^T$.
\end{enumerate}

\medskip
\noindent
La formulación precisa aquí las cotas globales de las derivadas de los
coeficientes en divergencia, que se utilizan al pasar a la forma no
divergente; no basta una cota puntual de $a^{ij}$, $b^i$, $c^i$ y $d$
para controlar automáticamente $\widetilde b$. Bajo $\widetilde c\le0$,
estas condiciones permiten aplicar el resultado clásico de existencia y
el principio del máximo \cite{gilbarg2001elliptic}. La constante de
estabilidad depende de $U$, $\lambda_0$ y una cota de $\widetilde b$;
esta última puede involucrar derivadas de $a^{ij}$ que no aparecen en $\nu$.

\section{Estabilidad uniforme para Dirichlet}
\label{sec:estabilidad_dirichlet}

\noindent\textbf{Procedencia y alcance.}
La estimación siguiente corresponde al principio de estabilidad usado por
\cite[Apéndice~D, Lema~D.1]{shin2020convergence}, donde se recurre al
principio del máximo y a una cota clásica de Gilbarg--Trudinger.
Aquí se reconstruye su demostración con una barrera explícita. Esta
reconstrucción permite indicar exactamente qué condiciones intervienen y de
qué cantidades depende la constante; no constituye un contraejemplo al
resultado elíptico original.

\medskip
\noindent\textbf{Enunciado (estabilidad $C^0$ para Dirichlet).}
Sea $U\subset\mathbb R^d$ un dominio acotado y supóngase que el operador
\emph{en forma no divergente}, obtenido mediante la expansión del Capítulo~I,
es
\[
\mathcal L[v]=\sum_{i,j=1}^d a^{ij}D_{ij}v+
\sum_{j=1}^d\widetilde b^{\,j}D_jv+\widetilde c\,v,
\qquad
\widetilde b^{\,j}=\sum_{i=1}^dD_i a^{ij}+b^j+c^j,
\quad
\widetilde c=\sum_{i=1}^dD_i b^i+d.
\]
Supongamos que $A=(a^{ij})$ es uniformemente elíptica, con constante
$\lambda_0>0$, que sus coeficientes y $\widetilde b$ son acotados y
continuos en $U$, y que $\widetilde c\leq 0$ es acotada y continua.
Las condiciones pertinentes se incluyen en la Hipótesis~3 precisada abajo.
Entonces existe $C_0>0$, independiente de $v$, tal que para
$v\in C^2(U)\cap C^0(\overline U)$ con $\mathcal L[v]$ acotada,
\begin{equation}\label{eq:estabilidad_error_general}
\|v\|_{C^0(\overline U)}
\leq \|v\|_{C^0(\partial U)}+
C_0\|\mathcal L[v]\|_{L^\infty(U)}.
\end{equation}
En particular, si $u^*$ es la solución clásica del problema de Dirichlet
$\mathcal L[u^*]=f$, $u^*|_{\partial U}=g$, y
$h\in C^2(U)\cap C^0(\overline U)$, entonces
\begin{equation}\label{eq:stability_C0}
\|u^*-h\|_{C^0(\overline U)}
\leq
\|g-h\|_{C^0(\partial U)}+
C_0\|f-\mathcal L[h]\|_{L^\infty(U)}.
\end{equation}
La constante puede elegirse en términos de $\lambda_0$, de la extensión
del dominio en una dirección y de una cota de $\widetilde b$, sin depender
del número de muestras ni de los parámetros neuronales.

\medskip
\noindent\textbf{Demostración mediante una barrera explícita.}
Sean $\ell=\inf_{x\in U}x_1$ y $R=1+\sup_{x\in U}x_1$.
Definamos
\[
B:=\sup_{x\in U}|\widetilde b^{\,1}(x)|<\infty,
\qquad
\gamma>\max\{1,2B/\lambda_0\},
\qquad
q(x):=1-\exp\bigl(\gamma(x_1-R)\bigr).
\]
Como $x_1\le R-1$ en $U$, se tiene $q\ge0$. Además,
$D_1q=-\gamma e^{\gamma(x_1-R)}$,
$D_{11}q=-\gamma^2 e^{\gamma(x_1-R)}$, y las demás derivadas
primeras y segundas se anulan. La elipticidad $a^{11}\ge\lambda_0$
y la condición $\widetilde c\le0$ proporcionan
\[
\begin{aligned}
\mathcal L[q]
&=-\gamma^2a^{11}e^{\gamma(x_1-R)}
-\gamma\widetilde b^{\,1}e^{\gamma(x_1-R)}
+\widetilde c\,q\\
&\le-e^{\gamma(x_1-R)}(\lambda_0\gamma^2-B\gamma)
\le-\delta,
\end{aligned}
\]
donde
\[
\delta:=e^{\gamma(\ell-R)}(\lambda_0\gamma^2-B\gamma)>0.
\]
Por lo tanto, $\psi=q/\delta$ cumple $\psi\ge0$,
$\mathcal L[\psi]\le-1$ y $\|\psi\|_\infty\le\delta^{-1}$.

Para $v$ como en el enunciado pongamos
$M:=\|v\|_{C^0(\partial U)}$ y
$F:=\|\mathcal L[v]\|_{L^\infty(U)}$.
Las funciones $w_\pm:=\pm v-M-F\psi$ son no positivas en
$\partial U$ y verifican
\[
\mathcal L[w_\pm]
=\pm\mathcal L[v]-M\widetilde c-F\mathcal L[\psi]
\ge -F+F=0
\quad\text{en }U.
\]
El principio del máximo débil para operadores uniformemente elípticos
con $\widetilde c\le0$ da $w_\pm\le0$ en $\overline U$.
Por consiguiente,
\[
|v(x)|\le M+F\psi(x),\qquad
\|v\|_{C^0(\overline U)}\le M+\delta^{-1}F.
\]
Esto prueba \eqref{eq:estabilidad_error_general} con
$C_0=\delta^{-1}$. Sustituyendo $v=u^*-h$ y utilizando
$\mathcal L[u^*]=f$ se obtiene \eqref{eq:stability_C0}.
\hfill$\square$

\medskip
\noindent\textbf{Observación (separación entre existencia y estabilidad).}
La condición de esfera exterior se emplea para obtener existencia y
continuidad hasta la frontera de una solución clásica con dato continuo,
según el Teorema~6.13 de Gilbarg--Trudinger. En cambio, una vez dadas
$u^*$ y $h$ con la regularidad indicada, la desigualdad de estabilidad
anterior solo utiliza que $U$ es acotado, la elipticidad uniforme,
$\widetilde b\in L^\infty(U)$ y $\widetilde c\le0$; no se necesitan
hipótesis de frontera $C^{2,\alpha}$. Esto permite aplicar la estimación
también en el cuadrado de los experimentos.

\section{Convergencia uniforme de los aproximantes (Teorema 3)}

\noindent\textbf{Enunciado (convergencia $C^0$ de la solución aproximada).}
Bajo las Hipótesis~1,~2 y~3, tomando explícitamente el operador
de frontera de Dirichlet $\mathcal B[h]=h|_{\partial U}$, con
$\Gamma=\partial U$ y $\widetilde c\le0$ en $U$, sea $h_m$ una sucesión de minimizadores
globales de las pérdidas regularizadas del Teorema~2, con los pesos allí
fijados y a lo largo de su sucesión de tamaños muestrales. Entonces,
con probabilidad uno respecto del muestreo i.i.d.,
\begin{equation}\label{eq:teorema3_convergencia_C0}
\|h_m-u^*\|_{C^0(\overline U)}\longrightarrow0.
\end{equation}
Para $d\ge2$ y $0<\alpha<1$, en el mismo evento de probabilidad uno,
la demostración proporciona además la cota conservadora
\begin{equation}\label{eq:teorema3_tasa_C0}
\|h_m-u^*\|_{C^0(\overline U)}
=O\!\left(m_r^{-\alpha^2/[d(2\alpha+d)]}\right).
\end{equation}
Esta tasa describe el límite idealizado de minimizadores globales bajo
las hipótesis de interpolación exacta y regularización del Teorema~2,
no la convergencia del algoritmo de entrenamiento.

\medskip
\noindent\textbf{Demostración.}
Sea $e_m=u^*-h_m$. La desigualdad \eqref{eq:stability_C0}, cuyo
factor $C_0$ depende solo del operador y del dominio, da
\[
\|e_m\|_{C^0(\overline U)}
\le
\|g-h_m\|_{C^0(\partial U)}+
C_0\|f-\mathcal L[h_m]\|_{C^0(\overline U)}.
\]
Por el Teorema~2 ambos sumandos convergen a cero en un único evento
de probabilidad uno, lo que prueba \eqref{eq:teorema3_convergencia_C0}.
Para $d\ge2$, sustituimos las estimaciones cuantitativas
\eqref{eq:shin_sup_rates}. Como
\[
\frac{\alpha^2}{d(2\alpha+d-1)}>
\frac{\alpha^2}{d(2\alpha+d)},
\]
la cota correspondiente a la frontera decae al menos tan rápidamente
como la del residuo interior. Obtenemos así
\eqref{eq:teorema3_tasa_C0}.
\hfill$\square$

\medskip
\noindent\textbf{Observación (procedencia de la demostración).}
La implicación ``convergencia uniforme de residuos y de frontera $\Rightarrow$
convergencia uniforme de aproximantes'' es la del Teorema~3.3 de
\cite{shin2020convergence}. La prueba precedente reconstruye de manera
autónoma el paso elíptico y explicita una cota de error. La reparación
referida a la comparabilidad de los pesos pertenece al Teorema~2 y se
transmite aquí a través de sus hipótesis; no es necesario atribuir un
error distinto al argumento de estabilidad de Shin.

\section{Convergencia en norma de energía con Dirichlet fuerte (Teorema 4)}
\label{sec:teorema4_dirichlet_fuerte}

El Teorema~3 relaciona el residuo y el error en norma uniforme mediante el
principio del máximo. Para el problema de Poisson con condiciones de Dirichlet
impuestas exactamente se dispone, además, de una estimación energética en
$H^1(U)$. Formulamos primero la parte determinista, que no requiere ni
muestreo aleatorio ni minimización de una pérdida neuronal.

\medskip
\noindent\textbf{Lema (estabilidad $H^1$ para Poisson con traza nula).}
Sea $U\subset\mathbb R^d$ un dominio acotado Lipschitz y sea $C_P>0$ una
constante de Poincaré para $H_0^1(U)$, de modo que
$\|v\|_{L^2(U)}\le C_P\|\nabla v\|_{L^2(U)}$ para todo
$v\in H_0^1(U)$. Si $e\in H_0^1(U)$ y su laplaciano distribucional
$\Delta e\in L^2(U)$, entonces
\begin{equation}\label{eq:teorema4_estabilidad_energia}
\boxed{\quad
\|e\|_{H^1(U)}\le C_P\sqrt{1+C_P^2}\,
\|\Delta e\|_{L^2(U)}.\quad}
\end{equation}
\noindent\emph{Demostración del lema.}
La identidad distribucional y la densidad de $C_c^\infty(U)$ en
$H_0^1(U)$ implican que, para toda $\varphi\in H_0^1(U)$,
\[
\int_U\nabla e\cdot\nabla\varphi\,dx
=-\int_U(\Delta e)\,\varphi\,dx.
\]
Tomando $\varphi=e$ y aplicando Cauchy--Schwarz y Poincaré,
\[
\|\nabla e\|_{L^2(U)}^2
\le\|\Delta e\|_{L^2(U)}\|e\|_{L^2(U)}
\le C_P\|\Delta e\|_{L^2(U)}\|\nabla e\|_{L^2(U)}.
\]
Si $\|\nabla e\|_2=0$, entonces $e=0$ por Poincaré. En otro caso
podemos dividir, obteniendo
$\|\nabla e\|_2\le C_P\|\Delta e\|_2$ y
$\|e\|_2\le C_P^2\|\Delta e\|_2$. La combinación cuadrática de
estas dos desigualdades proporciona
\eqref{eq:teorema4_estabilidad_energia}.\hfill$\square$

\medskip
\noindent\textbf{Enunciado (Teorema 4: convergencia $H^1$ con Dirichlet fuerte).}
Considérese $\mathcal L=\Delta$ en un dominio acotado Lipschitz
$U\subset\mathbb R^d$, con $\Gamma=\partial U$. Supóngase que
$u^*\in H^1(U)\cap C^0(\overline U)$ resuelve
$\Delta u^*=f$ en el sentido distribucional, con
$f\in L^2(U)$, y satisface $u^*|_{\partial U}=g$.
Se escogen las clases $\mathcal H_m$ del Teorema~2 dentro del espacio
afín de funciones con traza $g$ en toda $\partial U$; se exige que
dichas clases sigan satisfaciendo la Hipótesis~2, incluidos los
comparadores de pérdida empírica nula con regularidad uniforme.
Sea $h_m$ una elección mensurable de minimizadores globales de la
pérdida del Teorema~2 \emph{sobre esas clases con Dirichlet fuerte}.
Supóngase además que
$h_m\in H^1(U)\cap C^0(\overline U)$,
$\Delta h_m\in L^2(U)$; por construcción, cada aproximante satisface
\emph{exactamente} $h_m=g$ sobre $\partial U$.
Entonces, con probabilidad uno respecto del muestreo i.i.d. del
Teorema~2,
\begin{equation}\label{eq:teorema4_convergencia_h1}
\|h_m-u^*\|_{H^1(U)}\longrightarrow0.
\end{equation}
En particular, para $d\ge2$ y $0<\alpha<1$, a lo largo de la
sucesión de tamaños muestrales especificada en el Teorema~2, se obtiene
la tasa condicional
\begin{equation}\label{eq:teorema4_tasa_h1_general}
\|h_m-u^*\|_{H^1(U)}
=O\!\left(m_r^{-\alpha^2/[d(2\alpha+d)]}\right).
\end{equation}
De hecho, la Hipótesis~1 ya impone una cota inferior positiva
$\rho_r(x)\ge\rho_0>0$ casi en todas partes de $U$; por tanto,
en el marco de este teorema se dispone directamente de la tasa más favorable:
\begin{equation}\label{eq:teorema4_tasa_h1_densidad}
\|h_m-u^*\|_{H^1(U)}=O\!\left(m_r^{-\alpha/(2d)}\right).
\end{equation}
La cota de densidad se cumple, por ejemplo, para la medida
uniforme de probabilidad en $U$, con $\rho_0=|U|^{-1}$.
La tasa anterior es condicional al conjunto de hipótesis del
Teorema~2 y a la elección de clases con Dirichlet fuerte; no se
sigue de la mera parametrización de frontera.

\medskip
\noindent\textbf{Demostración.}
En un dominio Lipschitz existe un operador de traza continuo
$\operatorname{Tr}:H^1(U)\to H^{1/2}(\partial U)$ y se tiene
$H_0^1(U)=\ker(\operatorname{Tr})$. Puesto que $u^*$ y $h_m$
son continuas hasta la frontera y sus valores de contorno coinciden
con $g$, sus trazas también coinciden. Por tanto,
\[
e_m:=u^*-h_m\in H_0^1(U),\qquad
\Delta e_m=f-\Delta h_m=:r_m\in L^2(U).
\]
El lema anterior proporciona, para todo $m$,
\begin{equation}\label{eq:teorema4_residuo_h1}
\|h_m-u^*\|_{H^1(U)}
\le C_P\sqrt{1+C_P^2}\,\|f-\Delta h_m\|_{L^2(U)}.
\end{equation}
Por el Teorema~2, en un evento de probabilidad uno,
$\|r_m\|_{C^0(\overline U)}\to0$. Puesto que $U$ es acotado,
\[
\|r_m\|_{L^2(U)}\le |U|^{1/2}\|r_m\|_{C^0(\overline U)}\to0.
\]
Aplicando \eqref{eq:teorema4_residuo_h1} se deduce
\eqref{eq:teorema4_convergencia_h1}; la primera tasa se obtiene
sustituyendo la cota interior de \eqref{eq:shin_sup_rates}.

Para la tasa bajo la cota de densidad ya exigida por la Hipótesis~1, usamos la
convergencia de la \emph{pérdida poblacional}, sin pasar por la norma
uniforme. Por la condición $\rho_r\ge\rho_0$ y la positividad del
peso residual fijo $\lambda_r$,
\[
\begin{aligned}
\|r_m\|_{L^2(U)}^2
&=\int_U|r_m|^2\,dx
\le\rho_0^{-1}\int_U|r_m|^2\,d\mu_r\\
&\le(\rho_0\lambda_r)^{-1}
\operatorname{Loss}^{\mathrm{PINN}}(h_m;\lambda)
=O\!\left(m_r^{-\alpha/d}\right),
\end{aligned}
\]
donde la última igualdad procede de
\eqref{eq:shin_expected_rate}. Tomando raíces cuadradas y
sustituyendo en \eqref{eq:teorema4_residuo_h1} se demuestra
\eqref{eq:teorema4_tasa_h1_densidad}.
Las dos tasas son afirmaciones sobre minimizadores globales bajo las
hipótesis del Teorema~2, no sobre el número de iteraciones de Adam o
L-BFGS.\hfill$\square$

\medskip
\noindent\textbf{Observación (alcance de la regularidad y de Dirichlet fuerte).}
La regularidad $u^*\in C^2(U)\cap C^0(\overline U)$, por sí sola,
no implica $u^*\in H^1(U)$. Por ello se incluye explícitamente
esta última condición. Asimismo, la igualdad de Dirichlet debe ser
exacta en el sentido de trazas; una penalización de frontera con peso
finito no basta para deducir que $e_m\in H_0^1(U)$.
La arquitectura $h_\theta(x,y)=(1-x^2)(1-y^2)N_\theta(x,y)$,
utilizada en el Capítulo~IV para $g=0$ en $(-1,1)^2$, sí
impone Dirichlet fuerte. Sin embargo, esto no verifica por sí solo
la Hipótesis~2 de interpolación exacta ni demuestra que Adam o
L-BFGS produzcan minimizadores globales.

\medskip
\noindent\textbf{Observación (procedencia y comparación con Shin).}
El Teorema~3.4 y el Apéndice~E de \cite{shin2020convergence}
consideran un operador elíptico más general. En la prueba publicada
se utiliza una desigualdad energética tras integrar por partes.
Para justificar ese paso hay que disponer de regularidad $H^1$
de las funciones involucradas, hipótesis que aquí se hace expresa.
Nuestra proposición se restringe al caso Poisson y demuestra
la desigualdad de estabilidad \eqref{eq:teorema4_residuo_h1}
directamente para el error con traza nula. Las tasas
\eqref{eq:teorema4_tasa_h1_general} y
\eqref{eq:teorema4_tasa_h1_densidad} son consecuencias
cuantitativas de los resultados previamente establecidos en esta
tesis; no afirman falsedad del teorema general de Shin.
La segunda tasa utiliza específicamente la cota inferior de la
densidad de la medida poblacional, que ya se incluye en nuestra
Hipótesis~1; no se deduciría de una hipótesis de cobertura puramente
geométrica que omitiera esa cota.

\medskip
\noindent\textbf{Observación (convención de signo).}
El capítulo teórico utiliza $\mathcal L=\Delta$.
La versión equivalente $-\Delta u=f_{\mathrm{num}}$ se obtiene
poniendo $f=-f_{\mathrm{num}}$; ambas conducen a la misma
estimación \eqref{eq:teorema4_residuo_h1}, porque el signo no
modifica la norma del residuo.
